% ---------------------------------------------------------------------------
% Lecture 6: Stochastic Processes. Built on one picture: the table of sample
% paths (rows = outcomes omega, columns = times t), and for Markov chains the
% transition matrix as a table whose rows are conditional distributions (the
% row picture of the Conditional Expectation deck). Homework: quiz03.tex.
% Every number was checked
% with python3 (sympy / fractions / numpy); the script is quoted in comments
% where the number appears.
% ---------------------------------------------------------------------------
\newcommand{\Z}{\mathbb{Z}}
\newcommand{\N}{\mathbb{N}}
\newcommand{\Var}{\operatorname{Var}}
\newcommand{\Cov}{\operatorname{Cov}}
\newcommand{\PP}{\mathbf{P}}   % transition matrix
\newcommand{\st}[1]{\mathsf{#1}} % named states
% Fixed random-walk path (numpy default_rng(11), 400 fair +-1 steps), scaled as
% (k/n, S_k/sqrt(n)) for n=20 and n=400.
\newcommand{\walkTwenty}{(0.0000,0.000) (0.0500,-0.224) (0.1000,-0.447) (0.1500,-0.224) (0.2000,-0.447) (0.2500,-0.224) (0.3000,0.000) (0.3500,0.224) (0.4000,0.000) (0.4500,-0.224) (0.5000,-0.447) (0.5500,-0.671) (0.6000,-0.447) (0.6500,-0.224) (0.7000,-0.447) (0.7500,-0.224) (0.8000,-0.447) (0.8500,-0.224) (0.9000,0.000) (0.9500,0.224) (1.0000,0.447)}
\newcommand{\walkFourHundred}{(0.0000,0.000) (0.0025,-0.050) (0.0050,-0.100) (0.0075,-0.050) (0.0100,-0.100) (0.0125,-0.050) (0.0150,0.000) (0.0175,0.050) (0.0200,0.000) (0.0225,-0.050) (0.0250,-0.100) (0.0275,-0.150) (0.0300,-0.100) (0.0325,-0.050) (0.0350,-0.100) (0.0375,-0.050) (0.0400,-0.100) (0.0425,-0.050) (0.0450,0.000) (0.0475,0.050) (0.0500,0.100) (0.0525,0.150) (0.0550,0.100) (0.0575,0.050) (0.0600,0.100) (0.0625,0.050) (0.0650,0.100) (0.0675,0.150) (0.0700,0.100) (0.0725,0.150) (0.0750,0.100) (0.0775,0.050) (0.0800,0.100) (0.0825,0.050) (0.0850,0.100) (0.0875,0.050) (0.0900,0.100) (0.0925,0.150) (0.0950,0.200) (0.0975,0.250) (0.1000,0.300) (0.1025,0.350) (0.1050,0.400) (0.1075,0.350) (0.1100,0.300) (0.1125,0.250) (0.1150,0.300) (0.1175,0.350) (0.1200,0.300) (0.1225,0.350) (0.1250,0.300) (0.1275,0.350) (0.1300,0.400) (0.1325,0.450) (0.1350,0.400) (0.1375,0.450) (0.1400,0.500) (0.1425,0.550) (0.1450,0.600) (0.1475,0.650) (0.1500,0.600) (0.1525,0.650) (0.1550,0.600) (0.1575,0.650) (0.1600,0.600) (0.1625,0.550) (0.1650,0.500) (0.1675,0.550) (0.1700,0.600) (0.1725,0.550) (0.1750,0.500) (0.1775,0.450) (0.1800,0.400) (0.1825,0.450) (0.1850,0.500) (0.1875,0.450) (0.1900,0.400) (0.1925,0.450) (0.1950,0.500) (0.1975,0.550) (0.2000,0.500) (0.2025,0.550) (0.2050,0.500) (0.2075,0.450) (0.2100,0.400) (0.2125,0.350) (0.2150,0.300) (0.2175,0.350) (0.2200,0.300) (0.2225,0.250) (0.2250,0.300) (0.2275,0.350) (0.2300,0.400) (0.2325,0.450) (0.2350,0.400) (0.2375,0.450) (0.2400,0.400) (0.2425,0.350) (0.2450,0.300) (0.2475,0.250) (0.2500,0.300) (0.2525,0.250) (0.2550,0.200) (0.2575,0.250) (0.2600,0.300) (0.2625,0.350) (0.2650,0.300) (0.2675,0.250) (0.2700,0.200) (0.2725,0.150) (0.2750,0.200) (0.2775,0.150) (0.2800,0.200) (0.2825,0.150) (0.2850,0.100) (0.2875,0.150) (0.2900,0.100) (0.2925,0.050) (0.2950,0.000) (0.2975,0.050) (0.3000,0.000) (0.3025,-0.050) (0.3050,-0.100) (0.3075,-0.150) (0.3100,-0.100) (0.3125,-0.150) (0.3150,-0.200) (0.3175,-0.250) (0.3200,-0.300) (0.3225,-0.250) (0.3250,-0.200) (0.3275,-0.150) (0.3300,-0.200) (0.3325,-0.150) (0.3350,-0.200) (0.3375,-0.250) (0.3400,-0.200) (0.3425,-0.150) (0.3450,-0.100) (0.3475,-0.050) (0.3500,-0.100) (0.3525,-0.150) (0.3550,-0.100) (0.3575,-0.050) (0.3600,-0.100) (0.3625,-0.050) (0.3650,-0.100) (0.3675,-0.050) (0.3700,-0.100) (0.3725,-0.150) (0.3750,-0.200) (0.3775,-0.150) (0.3800,-0.100) (0.3825,-0.050) (0.3850,-0.100) (0.3875,-0.150) (0.3900,-0.100) (0.3925,-0.050) (0.3950,-0.100) (0.3975,-0.150) (0.4000,-0.200) (0.4025,-0.250) (0.4050,-0.300) (0.4075,-0.350) (0.4100,-0.300) (0.4125,-0.350) (0.4150,-0.300) (0.4175,-0.250) (0.4200,-0.300) (0.4225,-0.250) (0.4250,-0.300) (0.4275,-0.250) (0.4300,-0.300) (0.4325,-0.250) (0.4350,-0.200) (0.4375,-0.150) (0.4400,-0.200) (0.4425,-0.250) (0.4450,-0.300) (0.4475,-0.250) (0.4500,-0.300) (0.4525,-0.250) (0.4550,-0.300) (0.4575,-0.350) (0.4600,-0.300) (0.4625,-0.350) (0.4650,-0.300) (0.4675,-0.350) (0.4700,-0.300) (0.4725,-0.350) (0.4750,-0.300) (0.4775,-0.250) (0.4800,-0.200) (0.4825,-0.150) (0.4850,-0.100) (0.4875,-0.050) (0.4900,0.000) (0.4925,-0.050) (0.4950,-0.100) (0.4975,-0.050) (0.5000,0.000) (0.5025,-0.050) (0.5050,-0.100) (0.5075,-0.050) (0.5100,-0.100) (0.5125,-0.150) (0.5150,-0.100) (0.5175,-0.050) (0.5200,0.000) (0.5225,-0.050) (0.5250,0.000) (0.5275,0.050) (0.5300,0.000) (0.5325,-0.050) (0.5350,-0.100) (0.5375,-0.050) (0.5400,-0.100) (0.5425,-0.050) (0.5450,0.000) (0.5475,-0.050) (0.5500,-0.100) (0.5525,-0.050) (0.5550,-0.100) (0.5575,-0.050) (0.5600,-0.100) (0.5625,-0.050) (0.5650,-0.100) (0.5675,-0.050) (0.5700,0.000) (0.5725,-0.050) (0.5750,0.000) (0.5775,-0.050) (0.5800,-0.100) (0.5825,-0.150) (0.5850,-0.200) (0.5875,-0.150) (0.5900,-0.200) (0.5925,-0.250) (0.5950,-0.200) (0.5975,-0.250) (0.6000,-0.300) (0.6025,-0.350) (0.6050,-0.400) (0.6075,-0.350) (0.6100,-0.400) (0.6125,-0.350) (0.6150,-0.300) (0.6175,-0.350) (0.6200,-0.300) (0.6225,-0.350) (0.6250,-0.400) (0.6275,-0.450) (0.6300,-0.400) (0.6325,-0.450) (0.6350,-0.500) (0.6375,-0.550) (0.6400,-0.600) (0.6425,-0.650) (0.6450,-0.700) (0.6475,-0.650) (0.6500,-0.700) (0.6525,-0.750) (0.6550,-0.700) (0.6575,-0.650) (0.6600,-0.700) (0.6625,-0.750) (0.6650,-0.800) (0.6675,-0.750) (0.6700,-0.700) (0.6725,-0.650) (0.6750,-0.600) (0.6775,-0.650) (0.6800,-0.600) (0.6825,-0.650) (0.6850,-0.600) (0.6875,-0.550) (0.6900,-0.600) (0.6925,-0.550) (0.6950,-0.500) (0.6975,-0.450) (0.7000,-0.400) (0.7025,-0.450) (0.7050,-0.500) (0.7075,-0.450) (0.7100,-0.500) (0.7125,-0.450) (0.7150,-0.500) (0.7175,-0.550) (0.7200,-0.500) (0.7225,-0.450) (0.7250,-0.400) (0.7275,-0.350) (0.7300,-0.300) (0.7325,-0.250) (0.7350,-0.200) (0.7375,-0.150) (0.7400,-0.200) (0.7425,-0.250) (0.7450,-0.200) (0.7475,-0.250) (0.7500,-0.300) (0.7525,-0.350) (0.7550,-0.400) (0.7575,-0.450) (0.7600,-0.500) (0.7625,-0.550) (0.7650,-0.600) (0.7675,-0.550) (0.7700,-0.600) (0.7725,-0.550) (0.7750,-0.500) (0.7775,-0.450) (0.7800,-0.400) (0.7825,-0.450) (0.7850,-0.500) (0.7875,-0.450) (0.7900,-0.400) (0.7925,-0.450) (0.7950,-0.500) (0.7975,-0.450) (0.8000,-0.400) (0.8025,-0.450) (0.8050,-0.500) (0.8075,-0.550) (0.8100,-0.500) (0.8125,-0.450) (0.8150,-0.500) (0.8175,-0.450) (0.8200,-0.400) (0.8225,-0.350) (0.8250,-0.300) (0.8275,-0.250) (0.8300,-0.200) (0.8325,-0.250) (0.8350,-0.300) (0.8375,-0.350) (0.8400,-0.300) (0.8425,-0.350) (0.8450,-0.400) (0.8475,-0.350) (0.8500,-0.400) (0.8525,-0.450) (0.8550,-0.400) (0.8575,-0.350) (0.8600,-0.400) (0.8625,-0.450) (0.8650,-0.500) (0.8675,-0.450) (0.8700,-0.400) (0.8725,-0.350) (0.8750,-0.300) (0.8775,-0.350) (0.8800,-0.300) (0.8825,-0.250) (0.8850,-0.200) (0.8875,-0.150) (0.8900,-0.100) (0.8925,-0.150) (0.8950,-0.200) (0.8975,-0.250) (0.9000,-0.200) (0.9025,-0.150) (0.9050,-0.200) (0.9075,-0.250) (0.9100,-0.200) (0.9125,-0.250) (0.9150,-0.300) (0.9175,-0.350) (0.9200,-0.400) (0.9225,-0.350) (0.9250,-0.300) (0.9275,-0.350) (0.9300,-0.300) (0.9325,-0.250) (0.9350,-0.200) (0.9375,-0.150) (0.9400,-0.100) (0.9425,-0.150) (0.9450,-0.100) (0.9475,-0.150) (0.9500,-0.200) (0.9525,-0.150) (0.9550,-0.100) (0.9575,-0.150) (0.9600,-0.200) (0.9625,-0.250) (0.9650,-0.200) (0.9675,-0.150) (0.9700,-0.100) (0.9725,-0.150) (0.9750,-0.100) (0.9775,-0.050) (0.9800,0.000) (0.9825,-0.050) (0.9850,-0.100) (0.9875,-0.150) (0.9900,-0.100) (0.9925,-0.050) (0.9950,-0.100) (0.9975,-0.050) (1.0000,0.000)}

\begin{frame}{Stochastic Processes}
  \small
  Throughout, $(\Omega,\F,\Prob)$ is one fixed probability space. The lecture is organised around a single picture, a \textbf{table of sample paths}: one row for each outcome $\omega$, one column for each time $t$, and the value $X_t(\omega)$ in the cell.
  \bigskip
  \renewcommand{\arraystretch}{1.7}
  \begin{tabular}{@{}r@{\ \ }>{\raggedright\arraybackslash}p{4.0cm}>{\raggedright\arraybackslash}p{8.8cm}@{}}
    1. & \textbf{Stochastic process} & A family of random variables indexed by time. Rows are paths, columns are random variables. \\
    2. & \textbf{Random walk} & Sum of independent $\pm1$ steps. The reflection principle and first-step analysis. \\
    3. & \textbf{Brownian motion} & The limit of the random walk under $\sqrt n$ scaling; independent Gaussian increments. \\
    4. & \textbf{Risk of ruin} & The gambler's ruin problem, solved by first-step analysis. \\
    5. & \textbf{Markov chains} & The transition matrix is a table whose rows are conditional distributions. Powers, stationary distributions, classes, hitting times, convergence. \\
  \end{tabular}
\end{frame}

% ===========================================================================
\begin{frame}{Stochastic Process: Definition}
  \label{fr:procdef}
  \small
  \begin{block}{Definition}
    Let $T$ be a set of \textbf{times}, either $T=\{0,1,2,\dots\}$ (discrete time) or $T=[0,\infty)$ (continuous time). A \textbf{stochastic process} is a family
    \[
      (X_t)_{t\in T},\qquad X_t\colon\Omega\to\R\ \text{ a random variable for each } t\in T,
    \]
    all defined on the same $(\Omega,\F,\Prob)$. Equivalently it is one function of two arguments, $X\colon T\times\Omega\to\R$, $(t,\omega)\mapsto X_t(\omega)$.
  \end{block}
  \words{There is one random variable for each time, and all are defined on the same $\Omega$.}
  \begin{itemize}
    \item \textbf{For fixed $\omega$}, the function $T\to\R$, $t\mapsto X_t(\omega)$, is the \textbf{sample path} (or trajectory) of the outcome $\omega$: an ordinary sequence or an ordinary function of time.
    \item \textbf{For fixed $t$}, the function $\Omega\to\R$, $\omega\mapsto X_t(\omega)$, is the random variable $X_t$: what the process shows at time $t$.
    \item The \textbf{state space} $S$ is the set of values the $X_t$ take, so $X_t\colon\Omega\to S$: usually $S\subseteq\R$ ($S=\{-1,+1\}$, $\Z$, $\R$); for Markov chains a finite or countable set of named states.
    \item Shorthand, used throughout: for $A\subseteq S$, $\{X_t\in A\}$ is the event $\ev{X_t(\omega)\in A}$ and $\Prob(X_t\in A)=\Prob(\ev{X_t(\omega)\in A})$; $\Prob(X_s=x,\ X_t=y)$ is the probability of the intersection of the two such events.
  \end{itemize}
\end{frame}

% ---------------------------------------------------------------------------
\begin{frame}{The Central Picture: A Table of Sample Paths}
  \label{fr:table}
  \small
  \begin{columns}[T]
    \begin{column}{0.52\textwidth}
      Three fair coin tosses. $\Omega=\{\st H,\st T\}^3$, each $\omega$ having probability $\tfrac18$. Put $\xi_t\colon\Omega\to\{-1,+1\}$, $\xi_t(\omega)=+1$ if toss $t$ of $\omega$ is $\st H$ and $-1$ if $\st T$, and
      \[
        S_t=\xi_1+\dots+\xi_t,\qquad S_0=0,\qquad t\in\{0,1,2,3\}.
      \]
      \begin{itemize}
        \item A \textbf{row} is one outcome $\omega$ and shows the path $t\mapsto S_t(\omega)$.
        \item A \textbf{column} is one time $t$ and shows the random variable $S_t$: its entries, with the row probabilities, are the distribution of $S_t$.
        \item An event is a set of rows: $\{S_3=1\}$ is rows $2,3,5$, so $\Prob(S_3=1)=\tfrac38$.
      \end{itemize}
      \words{Properties of the paths are read along the rows of this table, and properties of the distribution at a fixed time are read down the columns.}
    \end{column}
    \begin{column}{0.46\textwidth}
      \centering
      \renewcommand{\arraystretch}{1.15}
      \begin{tabular}{@{}lc|cccc@{}}
        \toprule
        $\omega$ & $\Prob(\{\omega\})$ & $S_0$ & $S_1$ & $S_2$ & $S_3$ \\
        \midrule
        $\st{HHH}$ & $\tfrac18$ & $0$ & $1$ & $2$ & $3$ \\
        $\st{HHT}$ & $\tfrac18$ & $0$ & $1$ & $2$ & $1$ \\
        $\st{HTH}$ & $\tfrac18$ & $0$ & $1$ & $0$ & $1$ \\
        $\st{HTT}$ & $\tfrac18$ & $0$ & $1$ & $0$ & $-1$ \\
        $\st{THH}$ & $\tfrac18$ & $0$ & $-1$ & $0$ & $1$ \\
        $\st{THT}$ & $\tfrac18$ & $0$ & $-1$ & $0$ & $-1$ \\
        $\st{TTH}$ & $\tfrac18$ & $0$ & $-1$ & $-2$ & $-1$ \\
        $\st{TTT}$ & $\tfrac18$ & $0$ & $-1$ & $-2$ & $-3$ \\
        \bottomrule
      \end{tabular}
      \par\medskip
      {\footnotesize rows: paths $t\mapsto S_t(\omega)$; \quad columns: random variables $S_t$}
    \end{column}
  \end{columns}
\end{frame}

% ---------------------------------------------------------------------------
\begin{frame}{Reading the Table Two Ways, and What a Process Is Not}
  \label{fr:reading}
  \small
  \begin{columns}[T]
    \begin{column}{0.56\textwidth}
      \begin{block}{Row questions (about paths)}
        Whether the path reaches $2$, whether it returns to $0$, and what its maximum is. Each is a set of rows, hence an event: $\{\max_{t\le3}S_t\ge2\}$ is rows $1,2$, probability $\tfrac28$.
      \end{block}
      \begin{block}{Column questions (about one time)}
        The values of $\E[S_2]$ and $\Prob(S_3>0)$. Each is read from column $2$ or $3$ together with the row probabilities: $\E[S_2]=\tfrac18(2+2+0+0+0+0-2-2)=0$, $\Prob(S_3>0)=\tfrac48$.
      \end{block}
      \begin{block}{Mixed questions (about several columns)}
        $\Prob(S_1=1,\ S_3=1)$ is rows $2,3$: $\tfrac28$. Different columns are \emph{not} independent in general: $\Prob(S_1=1)\Prob(S_3=1)=\tfrac12\cdot\tfrac38\ne\tfrac28$.
      \end{block}
    \end{column}
    \begin{column}{0.42\textwidth}
      A stochastic process is \textbf{not}:
      \begin{itemize}
        \item a single random variable: it is one per time;
        \item a list of distributions, one per column: the distributions of the $S_t$ separately do not determine how columns move together, i.e.\ the joint distribution of the row;
        \item a fixed function of time: the path depends on $\omega$, and different rows are different functions.
      \end{itemize}
      \words{The process is the entire table; its columns are the random variables $X_t$ and its rows are the possible histories.}
    \end{column}
  \end{columns}
\end{frame}

% ---------------------------------------------------------------------------
\begin{frame}{The Distribution of a Process}
  \label{fr:fdd}
  \small
  \begin{block}{Definition}
    The \textbf{finite-dimensional distributions} of $(X_t)_{t\in T}$ are the joint distributions of $(X_{t_1},\dots,X_{t_n})$ for all $n$ and all times $t_1<\dots<t_n$ in $T$: in the discrete case the joint pmfs
    \[
      p_{X_{t_1},\dots,X_{t_n}}\colon\R^n\to[0,1],\qquad (x_1,\dots,x_n)\mapsto\Prob(X_{t_1}=x_1,\dots,X_{t_n}=x_n).
    \]
  \end{block}
  \words{The model specifies the joint distribution of every finite set of columns, and these joint distributions together determine every probability concerning the process.}
  \begin{block}{Definition (independence)}
    Random variables $X_1,\dots,X_n$ are \textbf{independent}, written $X_1\indep X_2\indep\dots\indep X_n$, if for all sets $A_1,\dots,A_n\subseteq\R$
    \[
      \Prob(X_1\in A_1,\dots,X_n\in A_n)=\Prob(X_1\in A_1)\cdots\Prob(X_n\in A_n).
    \]
    A sequence $(X_n)_{n\ge1}$ is independent if every finite sub-family is.
  \end{block}
  \begin{itemize}
    \item In the table, the coin tosses $\xi_1,\xi_2,\xi_3$ are independent, but the partial sums $S_1,S_2,S_3$ are not (previous slide). In practice a model is specified by independent building blocks together with a rule for combining them, rather than by listing all the joint pmfs.
  \end{itemize}
\end{frame}


% ===========================================================================
\begin{frame}{The Bernoulli Process and the Simple Random Walk}
  \label{fr:walk}
  \small
  \begin{block}{Definition}
    Fix $p\in(0,1)$ and put $q=1-p$. A \textbf{Bernoulli process} is an independent sequence $\xi_1,\xi_2,\dots$ of random variables $\xi_n\colon\Omega\to\{-1,+1\}$ with $\Prob(\xi_n=+1)=p$, $\Prob(\xi_n=-1)=q$ for every $n$. The \textbf{simple random walk} is the process $(S_n)_{n\ge0}$,
    \[
      S_n\colon\Omega\to\Z,\qquad S_0=0,\qquad S_n=\xi_1+\xi_2+\dots+\xi_n .
    \]
    It is \textbf{symmetric} (or fair) if $p=\tfrac12$.
  \end{block}
  \words{A coin is tossed repeatedly and a running score is kept, $+1$ for each head and $-1$ for each tail; $S_n$ is the score after $n$ tosses.}
  \begin{itemize}
    \item The table on slide~\ref{fr:table} is this process up to $n=3$. Each row is a path $n\mapsto S_n(\omega)$ that moves by $\pm1$ at each step.
    \item \textbf{Independent increments}: for $m<n$, $S_n-S_m=\xi_{m+1}+\dots+\xi_n$ depends only on tosses after $m$, so it is independent of $(S_0,\dots,S_m)$ and has the same distribution as $S_{n-m}$. After time $m$, the process $(S_{m+k}-S_m)_{k\ge0}$ is a simple random walk independent of $(S_0,\dots,S_m)$.
  \end{itemize}
\end{frame}

% ---------------------------------------------------------------------------
\begin{frame}{Distribution of $S_n$}
  \label{fr:distSn}
  \small
  \begin{block}{Theorem}
    Let $H_n\colon\Omega\to\{0,\dots,n\}$, $H_n=\sum_{k=1}^n\ind\{\xi_k=+1\}$, be the number of $+1$ steps among the first $n$: a sum of $n$ independent Bernoulli$(p)$ indicators, hence binomial$(n,p)$, and $S_n=H_n-(n-H_n)=2H_n-n$. Hence, for $k\in\{-n,-n+2,\dots,n\}$,
    \[
      \Prob(S_n=k)=\binom{n}{\tfrac{n+k}{2}}\,p^{(n+k)/2}q^{(n-k)/2},
      \qquad
      \E[S_n]=n(p-q),\qquad \Var(S_n)=4npq .
    \]
  \end{block}
  \words{$S_n$ is an affine function of a binomial random variable; its mean is of order $n$ and its standard deviation of order $\sqrt n$.}
  \begin{block}{Proof}
    $S_n=k$ iff $H_n=\tfrac{n+k}2$, which requires $n+k$ even. $\E[S_n]=2\E[H_n]-n=2np-n=n(p-q)$ and $\Var(S_n)=4\Var(H_n)=4npq$. \hfill$\square$
  \end{block}
  \begin{exampleblock}{Fair walk, $n=4$}
    % python3: Fraction(comb(4,(4+k)//2),16) for k=-4,-2,0,2,4 -> 1/16,1/4,3/8,1/4,1/16
    \begin{tabular}{@{}c|ccccc@{}}
      $k$ & $-4$ & $-2$ & $0$ & $2$ & $4$ \\ \hline
      $\Prob(S_4=k)$ & $\tfrac1{16}$ & $\tfrac4{16}$ & $\tfrac6{16}$ & $\tfrac4{16}$ & $\tfrac1{16}$
    \end{tabular}
    \qquad $\E[S_4]=0$, $\Var(S_4)=4$, standard deviation $2=\sqrt4$.
  \end{exampleblock}
\end{frame}

% ---------------------------------------------------------------------------
\begin{frame}{A Path Question: Reaching $3$ Within $10$ Steps}
  \label{fr:reflect}
  \small
  \begin{columns}[T]
    \begin{column}{0.58\textwidth}
      For the fair walk, consider $\Prob(\max_{k\le10}S_k\ge3)$. This is a row question, since the event depends on the entire path rather than on the value at a single time.
      \begin{block}{Direct computation}
        One can enumerate the $2^{10}=1024$ paths and count those whose running maximum reaches $3$. The following theorem replaces this count by a computation involving the distribution of $S_n$ only.
      \end{block}
      \begin{block}{Theorem (reflection principle, fair walk)}
        For integers $m\ge1$ and $n\ge1$,
        \[
          \Prob\Big(\max_{k\le n}S_k\ge m\Big)=2\,\Prob(S_n>m)+\Prob(S_n=m).
        \]
      \end{block}
      \words{The distribution of the maximum up to time $n$ is expressed through the distribution of $S_n$.}
    \end{column}
    \begin{column}{0.40\textwidth}
      \centering
      \begin{tikzpicture}[x=0.36cm,y=0.36cm,font=\scriptsize]
        \draw[soft!50] (0,3) -- (10,3) node[right,soft] {$m=3$};
        \draw[->] (0,-2.5) -- (0,5.5);
        \draw[->] (0,0) -- (10.6,0) node[below right] {$n$};
        \draw[accent,thick] (0,0)--(1,1)--(2,2)--(3,3)--(4,4)--(5,3)--(6,2)--(7,1)--(8,2)--(9,1)--(10,0);
        \draw[amber,thick,dashed] (3,3)--(4,2)--(5,3)--(6,4)--(7,5)--(8,4)--(9,5)--(10,6);
        \fill[accent] (3,3) circle (2pt);
        \node[left] at (0,3) {$3$};
        \node[left] at (0,0) {$0$};
        \node[above right,accent] at (10,0) {$S_{10}$};
        \node[right,amber] at (10,6) {$6-S_{10}$};
      \end{tikzpicture}
      \par\smallskip
      {\footnotesize After the first visit to $m$, reflect the rest of the path in the line $m$. This pairs each path that reaches $m$ and ends below $m$ with one ending above $m$, and the two have equal probability.}
    \end{column}
  \end{columns}
\end{frame}

% ---------------------------------------------------------------------------
\begin{frame}{Reflection Principle: Proof}
  \label{fr:reflectproof}
  \small
  \begin{block}{Proof}
    Let $A=\{\max_{k\le n}S_k\ge m\}$. Split $A$ by where the path ends:
    \[
      \Prob(A)=\Prob(A,\ S_n>m)+\Prob(A,\ S_n=m)+\Prob(A,\ S_n<m).
    \]
    In the first two terms the condition $A$ is redundant, since a path ending at or above $m$ has reached $m$; they equal $\Prob(S_n>m)$ and $\Prob(S_n=m)$.

    \smallskip
    For the third, let $\tau$ be the first time the path hits $m$. Map each path in $\{A,\ S_n<m\}$ to the path that agrees with it up to time $\tau$ and thereafter takes the opposite steps. The new path ends at $m+(m-S_n)>m$, and it has the same number of steps, each of probability $\tfrac12$, so the same probability $2^{-n}$. The map is a bijection onto $\{S_n>m\}$ (reverse it by reflecting after the first visit to $m$). Hence $\Prob(A,\ S_n<m)=\Prob(S_n>m)$. \hfill$\square$
  \end{block}
  \words{For the fair walk all $2^n$ paths are equally likely, so probabilities are proportional to counts, and reflection is a bijection between two sets of paths.}
  \begin{itemize}
    \item For $p\ne\tfrac12$ the reflected path has a different probability, and the identity fails; the argument then yields only inequalities.
    \item Brownian motion inherits the statement in the form $\Prob(\max_{s\le t}B_s\ge a)=2\Prob(B_t\ge a)$ (slide~\ref{fr:bmreflect}).
  \end{itemize}
\end{frame}

% ---------------------------------------------------------------------------
\begin{frame}{Return to the Problem: Reaching $3$ Within $10$ Steps}
  \label{fr:reflectreturn}
  \small
  \begin{exampleblock}{The answer}
    % python3: P(S_10>=4)=sum(comb(10,h),h=7..10)/1024=176/1024=11/64; P(S_10=3)=0 (parity); 2*11/64=11/32.
    % brute force over all 1024 paths: 352/1024 = 11/32.  Checked.
    $S_{10}$ is even, so $\Prob(S_{10}=3)=0$, and $\Prob(S_{10}>3)=\Prob(S_{10}\ge4)=\Prob(H_{10}\ge7)=\dfrac{120+45+10+1}{1024}=\dfrac{11}{64}$. Hence
    \[
      \Prob\Big(\max_{k\le10}S_k\ge3\Big)=2\cdot\tfrac{11}{64}+0=\tfrac{11}{32}\approx0.34 .
    \]
    Direct enumeration of the $1024$ paths gives $352$, in agreement.
  \end{exampleblock}
  \begin{block}{Corollary: the distribution of the maximum}
    With $M_n=\max_{k\le n}S_k$ and the theorem applied at $m$ and $m+1$,
    \[
      \Prob(M_n=m)=\Prob(M_n\ge m)-\Prob(M_n\ge m+1)=\Prob(S_n=m)+\Prob(S_n=m+1)\qquad(m\ge0).
    \]
  \end{block}
  \words{The distribution of $M_n$, a path quantity, is determined by the distribution of $S_n$, a single column.}
  \begin{itemize}
    \item Fair walk, $n=10$: $\Prob(M_{10}=0)=\Prob(S_{10}=0)+\Prob(S_{10}=1)=\tfrac{252}{1024}$, so about a quarter of all $10$-step paths never rise above their start.
  \end{itemize}
\end{frame}

% ---------------------------------------------------------------------------
\begin{frame}{First-Step Analysis: Expected Time to Leave $(-a,a)$}
  \label{fr:firststep}
  \small
  Fair walk started at $i\in\{-a+1,\dots,a-1\}$, position $X_n=i+S_n$. Write $\Prob_i(\cdot)$, $\E_i[\cdot]$ for probabilities and expectations when the start is $i$ (the same notation is used for Markov chains below). \textbf{Exit time}: $\tau\colon\Omega\to\{0,1,\dots\}\cup\{\infty\}$, $\tau=\min\{n\ge0 : |X_n|\ge a\}$ ($\infty$ if there is no such $n$); $m_i=\E_i[\tau]$.
  \begin{block}{Direct computation}
    $\Prob_i(\tau=n)$ is found by counting the paths that stay strictly inside $(-a,a)$ for $n-1$ steps and exit at step $n$; then $\sum_n n\,\Prob_i(\tau=n)$. These counts have no simple closed form.
  \end{block}
  \begin{block}{First-step analysis}
    Condition on the first step, as for the waiting time for $\st{HH}$ (Conditional Expectation lecture, Application 2). After one step the walk is at $i\pm1$ and, by independent increments, the remaining steps $\xi_2,\xi_3,\dots$ form a simple random walk independent of $\xi_1$, so the expected exit time from that point is $m_{i\pm1}$:
    \[
      m_i=1+\tfrac12 m_{i+1}+\tfrac12 m_{i-1}\quad(-a<i<a),\qquad m_{-a}=m_a=0 .
    \]
  \end{block}
  \words{The expected exit time from $i$ is one step plus the average of the expected exit times from the two neighbouring positions.}
  \begin{itemize}
    \item The unknowns $m_{-a+1},\dots,m_{a-1}$ appear on both sides, so this is a system of $2a-1$ linear equations rather than an explicit formula; every first-step argument in this lecture has this form.
  \end{itemize}
\end{frame}

% ---------------------------------------------------------------------------
\begin{frame}{First-Step Analysis: The Solution}
  \label{fr:exitsol}
  \small
  \begin{block}{Theorem}
    $m_i=\E_i[\tau]=a^2-i^2$ for $-a\le i\le a$; in particular $\E_0[\tau]=a^2$.
  \end{block}
  \begin{block}{Proof}
    Try $m_i=c-i^2$: then $1+\tfrac12\big(c-(i+1)^2\big)+\tfrac12\big(c-(i-1)^2\big)=1+c-i^2-1=c-i^2$, so the equation holds for every $c$, and $m_{\pm a}=0$ forces $c=a^2$.
    The system is $2a-1$ linear equations in $2a-1$ unknowns; it has a unique solution (uniqueness: if $m,m'$ both solve it, $d=m-m'$ satisfies $d_i=\tfrac12(d_{i+1}+d_{i-1})$ with $d_{\pm a}=0$, so $d$ is linear and vanishes). \hfill$\square$
  \end{block}
  \words{The expected exit time is a quadratic function of the starting point, and from the centre it equals $a^2$.}
  % python3 sympy: a=3 linear system -> m0=9, m±1=8, m±2=5.  Checked.
  \begin{exampleblock}{$a=3$}
    \begin{tabular}{@{}c|ccccccc@{}}
      $i$ & $-3$ & $-2$ & $-1$ & $0$ & $1$ & $2$ & $3$ \\ \hline
      $m_i$ & $0$ & $5$ & $8$ & $9$ & $8$ & $5$ & $0$
    \end{tabular}
    \par\smallskip
    From $0$ the walk takes $9$ steps on average to reach $\pm3$; from $\pm2$ it takes $5$.\\
    This is consistent with slide~\ref{fr:distSn}: the walk travels a distance of order $\sqrt n$ in $n$ steps, so a distance $a$ takes a number of steps of order $a^2$.
  \end{exampleblock}
\end{frame}


% ===========================================================================
\begin{frame}{From the Walk to Brownian Motion: The $\sqrt n$ Scaling}
  \label{fr:scaling}
  \small
  \begin{columns}[T]
    \begin{column}{0.56\textwidth}
      A particle takes a fair $\pm1$ step every $\tfrac1n$ seconds, each of size $\tfrac1{\sqrt n}$. We ask what its path looks like when $n$ is large, and with what probability it is above $0$ at time $t=1$ and again at time $t=2$.
      \begin{block}{Rescaling}
        Position at time $t$ is $B^{(n)}_t=\tfrac1{\sqrt n}S_{\lfloor nt\rfloor}$. By slide~\ref{fr:distSn}, $\E[B^{(n)}_t]=0$ and $\Var(B^{(n)}_t)=\tfrac{\lfloor nt\rfloor}{n}\to t$. The central limit theorem (CLT) states: for i.i.d.\ $\xi_k$ with mean $0$ and variance $1$, $\Prob(S_m/\sqrt m\le x)\to\Phi(x)$ for every $x$, with $\Phi\colon\R\to[0,1]$ the $N(0,1)$ cdf. Applied with $m=\lfloor nt\rfloor$ it gives $\Prob\big(B^{(n)}_t\le x\big)\to\Phi\big(x/\sqrt t\big)$: \ $B^{(n)}_t\to N(0,t)$ in distribution.
        Increments over disjoint time intervals are sums of disjoint blocks of steps, so they are independent, with variance equal to the length of the interval.
      \end{block}
      \words{Division by $\sqrt n$ is the scaling under which the variance at time $t$ has a finite non-zero limit; the limiting process has normal marginals and independent increments.}
    \end{column}
    \begin{column}{0.42\textwidth}
      \centering
      \begin{tikzpicture}[x=4.5cm,y=0.9cm,font=\scriptsize]
        \draw[->] (0,-1.3) -- (0,1.3) node[above] {$B^{(n)}_t$};
        \draw[->] (0,0) -- (1.05,0) node[right] {$t$};
        \draw[amber,thick] plot coordinates {\walkTwenty};
        \draw[accent,thick] plot coordinates {\walkFourHundred};
        \node[amber] at (0.86,0.95) {$n=20$};
        \node[accent] at (0.86,-1.05) {$n=400$};
        \node[below] at (1,0) {$1$};
        \node[left] at (0,1) {$1$};
        \node[left] at (0,-1) {$-1$};
      \end{tikzpicture}
      \par\smallskip
      {\footnotesize The same coin sequence, $20$ and $400$ steps, each rescaled to fit $[0,1]$ with steps of size $1/\sqrt n$. At time $1$ both have variance $1$.}
    \end{column}
  \end{columns}
\end{frame}

% ---------------------------------------------------------------------------
\begin{frame}{The Invariance Principle}
  \label{fr:invariance}
  \small
  \begin{block}{Theorem (invariance principle, Donsker 1951; informal statement)}
    Let $\xi_1,\xi_2,\dots$ be independent with $\E[\xi_k]=0$ and $\Var(\xi_k)=1$, $S_n=\xi_1+\dots+\xi_n$, and $B^{(n)}_t=\tfrac1{\sqrt n}S_{\lfloor nt\rfloor}$. Then, as $n\to\infty$, the random path $t\mapsto B^{(n)}_t$ converges in distribution to a random continuous path $t\mapsto B_t$, and for all times $t_1<\dots<t_k$ and sets $A_1,\dots,A_k$,
    \[
      \Prob\big(B^{(n)}_{t_1}\in A_1,\dots,B^{(n)}_{t_k}\in A_k\big)\ \to\ \Prob\big(B_{t_1}\in A_1,\dots,B_{t_k}\in A_k\big).
    \]
    The limit process $B$ is the same for every step distribution with mean $0$ and variance $1$.
  \end{block}
  \words{On a large scale, every random walk with mean-zero, unit-variance steps has approximately the same distribution of paths, and the limiting distribution is that of Brownian motion.}
  \begin{itemize}
    \item The theorem is not proved here. The name refers to the fact that the limit does not depend on the step distribution; the central limit theorem is the special case of a single time $t_1=1$.
    \item The limit inherits three properties from the walk: independent increments over disjoint intervals (sums over disjoint blocks of steps), $\Var(\text{increment})=\text{length}$, and a start at $0$. Since the step size $1/\sqrt n$ tends to $0$, the limiting paths are continuous. These are the four properties in the definition on the next slide.
  \end{itemize}
\end{frame}

% ---------------------------------------------------------------------------
\begin{frame}{Brownian Motion: Definition}
  \label{fr:bmdef}
  \small
  \begin{block}{Definition}
    A \textbf{(standard) Brownian motion} is a process $B\colon[0,\infty)\times\Omega\to\R$, $(t,\omega)\mapsto B_t(\omega)$, such that
    \begin{enumerate}
      \item[(i)] $B_0=0$;
      \item[(ii)] \textbf{independent increments}: for all $0\le t_0<t_1<\dots<t_n$, the random variables $B_{t_1}-B_{t_0},\ B_{t_2}-B_{t_1},\dots,B_{t_n}-B_{t_{n-1}}$ are independent;
      \item[(iii)] \textbf{Gaussian increments}: for $0\le s<t$, $B_t-B_s\sim N(0,t-s)$, i.e.\ it has pdf $x\mapsto\frac{1}{\sqrt{2\pi(t-s)}}e^{-x^2/(2(t-s))}$;
      \item[(iv)] \textbf{continuous paths}: for every $\omega$, the function $[0,\infty)\to\R$, $t\mapsto B_t(\omega)$, is continuous.
    \end{enumerate}
  \end{block}
  \words{Over any time interval the displacement is normal with variance equal to the length of the interval, displacements over disjoint intervals are independent, and the paths have no jumps.}
  \begin{itemize}
    \item Existence is a theorem (Wiener, 1923); we take it as given. The four properties determine all finite-dimensional distributions (slide~\ref{fr:fdd}), hence every probability about $B$.
  \end{itemize}
\end{frame}

% ---------------------------------------------------------------------------
\begin{frame}{What $B$ Is: Reading the Table}
  \label{fr:bmreading}
  \small
  The table of paths now has a continuum of columns ($t\ge0$) and of rows ($\omega$); the two readings are unchanged.
  \begin{block}{Column $t$: the random variable $B_t\colon\Omega\to\R$}
    $B_t=B_t-B_0\sim N(0,t)$ by (i) and (iii). Its pdf is $f_{B_t}\colon\R\to[0,\infty)$, $f_{B_t}(x)=\frac1{\sqrt{2\pi t}}e^{-x^2/(2t)}$, so $\Prob(a\le B_t\le b)=\int_a^b f_{B_t}(x)\dd x$.
  \end{block}
  \begin{block}{Row $\omega$: the path $[0,\infty)\to\R$, $t\mapsto B_t(\omega)$}
    A continuous function starting at $0$. Different rows are different functions; a typical one is as irregular as the $n=400$ picture at every magnification. Row questions include the maximum up to time $t$, the first time the path reaches $a$, and the time spent above $0$.
  \end{block}
  \begin{block}{The two arguments $t$ and $\omega$}
    In $B_t(\omega)$, $t$ is the time (which column) and $\omega$ is the outcome (which row). For fixed $\omega$, $t\mapsto B_t(\omega)$ is a function of time; for fixed $t$, $\omega\mapsto B_t(\omega)$ is a random variable. An increment $B_t-B_s$ is a difference of two columns.
  \end{block}
  \words{The properties below all follow from writing values of $B$ as sums of independent increments.}
\end{frame}

% ---------------------------------------------------------------------------
\begin{frame}{What $B$ Is Not}
  \label{fr:bmnot}
  \small
  \begin{columns}[T]
    \begin{column}{0.40\textwidth}
      Brownian motion is \textbf{not}:
      \begin{itemize}
        \item a process with independent \emph{values}: $B_s$ and $B_t$ are strongly dependent (next slide); it is the \emph{increments} over disjoint intervals that are independent;
        \item stationary in value: $\Var(B_t)=t$ grows; it is the increments whose distribution depends only on the interval length;
        \item smooth: no path has a derivative anywhere (slide~\ref{fr:bmrough});
        \item a single normal random variable: it is one normal per $t$, and jointly they form a Gaussian family.
      \end{itemize}
    \end{column}
    \begin{column}{0.58\textwidth}
      \footnotesize
      \renewcommand{\arraystretch}{1.3}
      \begin{tabular}{@{}ll>{\raggedright\arraybackslash}p{2.4cm}@{}}
        \toprule
        & walk $S_n$ & Brownian $B_t$ \\
        \midrule
        time set $T$ & $\{0,1,2,\dots\}$ & $[0,\infty)$ \\
        state space & $\Z$ & $\R$ \\
        start & $S_0=0$ & $B_0=0$ \\
        increment over $(s,t]$ & sum of $t-s$ steps & $N(0,t-s)$ \\
        $\Var$ of increment & $t-s$ (fair) & $t-s$ \\
        disjoint increments & independent & independent \\
        path & moves by $\pm1$ & continuous, nowhere differentiable \\
        \bottomrule
      \end{tabular}
    \end{column}
  \end{columns}
  \words{$B$ has the same increment structure as the walk, in continuous time and with continuous state space.}
\end{frame}


% ---------------------------------------------------------------------------
\begin{frame}{Basic Properties: Marginals and Covariance}
  \label{fr:bmcov}
  \small
  \begin{block}{Theorem}
    \begin{enumerate}
      \item[(a)] $B_t\sim N(0,t)$: $\E[B_t]=0$, $\Var(B_t)=t$.
      \item[(b)] $\Cov(B_s,B_t)=\E[B_sB_t]=\min(s,t)$ for all $s,t\ge0$.
    \end{enumerate}
  \end{block}
  \words{Each $B_t$ is $N(0,t)$, and the covariance of two values is the earlier of the two times.}
  \begin{block}{Proof}
    (a) is (i) and (iii) with $s=0$. (b) For $s<t$ write $B_t=B_s+(B_t-B_s)$. The increment is independent of $B_s$ (property (ii) applied to $0<s<t$) and has mean $0$, so
    \[
      \E[B_sB_t]=\E[B_s^2]+\E\big[B_s\,(B_t-B_s)\big]=s+\E[B_s]\,\E[B_t-B_s]=s+0=\min(s,t).
    \]
    Since $\E[B_s]=0$, $\Cov(B_s,B_t)=\E[B_sB_t]$. \hfill$\square$
  \end{block}
  \begin{itemize}
    \item The correlation of $B_s$ and $B_t$ is $\min(s,t)/\sqrt{st}=\sqrt{s/t}$ for $s<t$: $B_1$ and $B_2$ have correlation $1/\sqrt2\approx0.71$; $B_1$ and $B_{100}$ have correlation $0.1$. The values are dependent although the increments are independent.
    \item Alternatively, $\E[B_sB_t]$ is a two-dimensional integral against the joint pdf of the next slide; the proof above avoids it.
  \end{itemize}
\end{frame}

% ---------------------------------------------------------------------------
\begin{frame}{Basic Properties: Joint Probability Density Function (pdf) of Several Columns}
  \label{fr:bmjoint}
  \small
  \begin{block}{Theorem}
    For $0=t_0<t_1<\dots<t_n$, the vector $(B_{t_1},\dots,B_{t_n})$ has joint pdf $f\colon\R^n\to[0,\infty)$,
    \[
      f(x_1,\dots,x_n)=\prod_{k=1}^n\frac{1}{\sqrt{2\pi(t_k-t_{k-1})}}\exp\!\Big(-\frac{(x_k-x_{k-1})^2}{2(t_k-t_{k-1})}\Big),\qquad x_0=0 .
    \]
  \end{block}
  \words{Any finite collection of values of $B$ has a joint density, obtained from the product of the increment densities by the substitution $d_k=x_k-x_{k-1}$.}
  \begin{block}{Proof}
    The increments $D_k=B_{t_k}-B_{t_{k-1}}$, $k=1,\dots,n$, are independent $N(0,t_k-t_{k-1})$ by (ii) and (iii), so $(D_1,\dots,D_n)$ has the product pdf $\prod_k\frac{1}{\sqrt{2\pi(t_k-t_{k-1})}}e^{-d_k^2/(2(t_k-t_{k-1}))}$.
    The values are $B_{t_k}=D_1+\dots+D_k$: a linear map from $(d_1,\dots,d_n)$ to $(x_1,\dots,x_n)$ with inverse $d_k=x_k-x_{k-1}$, of determinant $1$. A determinant-one change of variables leaves a pdf unchanged apart from the substitution. \hfill$\square$
  \end{block}
\end{frame}

% ---------------------------------------------------------------------------
\begin{frame}{Basic Properties: Scaling, Symmetry, Restart}
  \label{fr:bmscaling}
  \small
  \begin{block}{Theorem}
    Let $B$ be a Brownian motion. Each of the following is again a Brownian motion:
    \begin{enumerate}
      \item[(a)] \textbf{Scaling}: $\big(\tfrac{1}{\sqrt c}\,B_{ct}\big)_{t\ge0}$ for any fixed $c>0$.
      \item[(b)] \textbf{Symmetry}: $(-B_t)_{t\ge0}$.
      \item[(c)] \textbf{Restart}: $(B_{s+t}-B_s)_{t\ge0}$ for any fixed $s\ge0$, and it is independent of $(B_u)_{0\le u\le s}$.
    \end{enumerate}
  \end{block}
  \words{Rescaling, reflecting in the time axis, and restarting from a fixed time each give a process with the same distribution as $B$.}
  \begin{block}{Proof}
    Check (i)--(iv). (a) $\tfrac1{\sqrt c}B_{c\cdot0}=0$; increments $\tfrac1{\sqrt c}(B_{ct}-B_{cs})$ over disjoint intervals are independent, normal with variance $\tfrac1c(ct-cs)=t-s$; paths are continuous. (b) $-N(0,\sigma^2)$ is $N(0,\sigma^2)$. (c) Increments of the new process are increments of $B$ over intervals inside $[s,\infty)$, independent of increments over intervals inside $[0,s]$, and $B_u$ for $u\le s$ is such an increment. \hfill$\square$
  \end{block}
  \smallskip
  By (a) with $c=n$, $\tfrac1{\sqrt n}B_{nt}$ is again Brownian: the scaling of slide~\ref{fr:scaling} applied to $B$ itself.
  (c) is the continuous-time form of the restart property used in first-step analysis: after any fixed time $s$, $B_{s+t}-B_s$ is a Brownian motion independent of the past.
\end{frame}

% ---------------------------------------------------------------------------
\begin{frame}{Return to the Problem: $\Prob(B_2>0\mid B_1=1)$}
  \label{fr:bmcond}
  \small
  \begin{block}{By the independent increment}
    $B_2=B_1+(B_2-B_1)$, and $B_2-B_1\sim N(0,1)$ is independent of $B_1$. On the row $\{B_1=1\}$, $B_2=1+(B_2-B_1)$, so given $B_1=1$, $B_2\sim N(1,1)$ and
    % python3: Phi(1)=0.8413447; sympy integral of conditional density agrees 0.841345
    \[
      \Prob(B_2>0\mid B_1=1)=\Prob\big(N(1,1)>0\big)=\Prob\big(N(0,1)>-1\big)=\Phi(1)\approx0.841,
    \]
    with $\Phi$ the $N(0,1)$ cdf of slide~\ref{fr:scaling}.
  \end{block}
  \begin{block}{Via the joint pdf}
    By slide~\ref{fr:bmjoint} with $t_1=1$, $t_2=2$, $(B_1,B_2)$ has joint pdf $f(x,y)=\frac1{2\pi}e^{-x^2/2}e^{-(y-x)^2/2}$, a bivariate normal with covariance matrix $\begin{pmatrix}1&1\\1&2\end{pmatrix}$. Form $f_{B_2\mid B_1}(y\mid x)=f(x,y)/f_{B_1}(x)=\frac1{\sqrt{2\pi}}e^{-(y-x)^2/2}$ and integrate over $y>0$ at $x=1$. The conditional density is that of $N(x,1)$, and the answer is the same.
  \end{block}
  \words{Conditioning on $B_1$ fixes the first increment and does not affect the distribution of the second; conditional probabilities for $B$ are computed by splitting values into increments.}
  \begin{itemize}
    \item The same split gives the conditional distribution of $B_t$ given $B_s=x$ for any $s<t$: $N(x,\,t-s)$. The conditional distribution depends on the past only through the present value.
  \end{itemize}
\end{frame}

% ---------------------------------------------------------------------------
\begin{frame}{Return to the Opening Problem: $\Prob(B_1>0,\ B_2>0)$}
  \label{fr:bmopen}
  \small
  Let $X=B_1$ and $Z=B_2-B_1$: independent, both $N(0,1)$, and $B_2=X+Z$. The event is $\{X>0,\ X+Z>0\}$.
  \begin{block}{By the independent increment and symmetry}
    Split according to the sign of $Z$: $\{X>0,\ X+Z>0\}=\{X>0,\ Z>0\}\cup\{X>0,\ Z<0,\ X>-Z\}$.
    The first event has probability $\tfrac12\cdot\tfrac12=\tfrac14$ by independence. For the second: on the row $\{X>0,\ Z<0\}$, which has probability $\tfrac14$, the pair $(X,-Z)$ consists of two independent positive copies of the same distribution, so $X>-Z$ has conditional probability $\tfrac12$ by symmetry. Hence
    % python3: 1/4 + asin(1/sqrt 2)/(2 pi) = 0.375 = 3/8.  Checked.
    \[
      \Prob(B_1>0,\ B_2>0)=\tfrac14+\tfrac14\cdot\tfrac12=\tfrac38 .
    \]
  \end{block}
  \begin{block}{Via the joint pdf}
    Alternatively, integrate the joint pdf $\frac1{2\pi}e^{-x^2/2}e^{-(y-x)^2/2}$ over the quadrant $x>0$, $y>0$. This is the orthant probability of a bivariate normal with correlation $\rho=1/\sqrt2$, which by a polar-coordinate computation equals $\tfrac14+\frac{\arcsin\rho}{2\pi}=\tfrac14+\tfrac{\pi/4}{2\pi}=\tfrac38$.
  \end{block}
  \words{Given that the particle is above $0$ at time $1$, it is above $0$ at time $2$ with probability $\tfrac{3/8}{1/2}=\tfrac34$ rather than $\tfrac12$; the values $B_1$ and $B_2$ are dependent although the increments are independent.}
\end{frame}

% ---------------------------------------------------------------------------
\begin{frame}{The Reflection Principle for Brownian Motion}
  \label{fr:bmreflect}
  \small
  \begin{block}{Theorem (reflection principle)}
    For $a>0$ and $t>0$: \quad $\Prob\big(\max_{s\le t}B_s\ge a\big)=2\,\Prob(B_t\ge a)=\Prob(|B_t|\ge a)$.
  \end{block}
  \words{The maximum up to time $t$ has the same distribution as $|B_t|$, a single column of the table.}
  \begin{itemize}
    \item This is the limit of slide~\ref{fr:reflect}; the term $\Prob(S_n=m)$ vanishes under the $\sqrt n$ scaling. Example: $\Prob(\max_{s\le1}B_s\ge1)=2(1-\Phi(1))\approx0.317$; a direct computation would involve a continuum of columns.
  \end{itemize}
  % python3: 2*Phi(a/sqrt t)-1 ~ 2*phi(0)*a/sqrt t = a*sqrt(2/(pi t)); integral over t diverges; 2(1-Phi(1)) = 0.3173
  \begin{block}{Corollary (maximum and first hitting time)}
    Let $M_t=\max_{s\le t}B_s$ and $T_a\colon\Omega\to[0,\infty]$, $T_a=\min\{s\ge0 : B_s=a\}$. Then $M_t$ has the distribution of $|B_t|$, and
    \[
      \Prob(T_a\le t)=\Prob(M_t\ge a)=2\big(1-\Phi(a/\sqrt t)\big)\ \xrightarrow{\ t\to\infty\ }\ 1,\qquad\text{yet}\qquad \E[T_a]=\int_0^\infty\Prob(T_a>t)\dd t=\infty ,
    \]
    since $\Prob(T_a>t)=2\Phi(a/\sqrt t)-1\approx a\sqrt{2/(\pi t)}$ for large $t$.
  \end{block}
  \words{Every level is reached with probability one, but the mean time to reach it is infinite, as for the fair walk on slide~\ref{fr:regimes}.}
\end{frame}

% ---------------------------------------------------------------------------
\begin{frame}{Why the Paths Are Nowhere Differentiable}
  \label{fr:bmrough}
  \small
  \begin{block}{Theorem (nowhere differentiable paths)}
    With probability one, $t\mapsto B_t(\omega)$ is nowhere differentiable (Paley, Wiener, Zygmund, 1933; not proved here).
  \end{block}
  \words{No path of $B$ has a derivative at any time.}
  \begin{block}{Why}
    By property (iii) of the definition (slide~\ref{fr:bmdef}), $\frac{B_{t+h}-B_t}{h}\sim N\big(0,\tfrac1h\big)$: the difference quotient has standard deviation $1/\sqrt h\to\infty$ as $h\to0$, so it cannot converge.
  \end{block}
\end{frame}

% ---------------------------------------------------------------------------
\begin{frame}{The Squared Increments Sum to $t$}
  \label{fr:bmqv}
  \small
  \begin{block}{Theorem}
    For $t>0$ and $\Delta_k=B_{kt/n}-B_{(k-1)t/n}$, $k=1,\dots,n$:
    \ $\E\big[\sum_{k=1}^n\Delta_k^2\big]=t$ and $\Var\big(\sum_{k=1}^n\Delta_k^2\big)=\dfrac{2t^2}{n}\to0$.
    % python3 comment: Var of (N(0,v))^2 = 2v^2, so var of the sum is n*2*(t/n)^2 = 2t^2/n
  \end{block}
  \begin{block}{Proof}
    The $\Delta_k$ are independent $N(0,\tfrac tn)$ with $\E[\Delta_k^2]=\tfrac tn$ and $\Var(\Delta_k^2)=2(\tfrac tn)^2$; sum over $k$. \hfill$\square$
  \end{block}
  \words{The squared increments of a Brownian path over $[0,t]$ sum to about $t$; for a differentiable function they tend to $0$ as $n\to\infty$.}
  % python3: E|N(0,t/n)| = sqrt(2t/(pi n)); n of them: sqrt(2tn/pi) -> inf. Walk: sum of squared steps = n exactly; scaled by 1/sqrt n: n*(1/n)=1.
  \begin{itemize}
    \item The $n$ steps of $S_n$ have squared increments $(\pm1)^2$ summing to $n$, so the rescaled walk $B^{(n)}$ of slide~\ref{fr:scaling} has $\sum(\Delta B^{(n)})^2=n\cdot\tfrac1n=1$ over $[0,1]$, for every $n$, and so does the limit.
    \item The increments themselves do not add up: $\E|\Delta_k|=\sqrt{2t/(\pi n)}$ (the mean of $|N(0,\sigma^2)|$ is $\sigma\sqrt{2/\pi}$), so $\E\big[\sum_{k=1}^n|\Delta_k|\big]=\sqrt{2tn/\pi}\to\infty$: over any interval the path has infinite length.
  \end{itemize}
\end{frame}


% ===========================================================================
\begin{frame}{Risk of Ruin: The Problem}
  \label{fr:ruinprob}
  \small
  A gambler with fortune $\pounds i$ bets $\pounds1$ per round, winning $\pounds1$ with probability $p$ and losing it with probability $q=1-p$, rounds independent. She stops when her fortune reaches $0$ (\textbf{ruin}) or a target $N>i$.
  \begin{exampleblock}{The opening question}
    Red at American roulette: $p=\tfrac{18}{38}$. She starts with $\pounds50$, aims for $\pounds100$, and stakes one pound per spin. We want the probability that she reaches the target, and the expected number of spins the game lasts.
  \end{exampleblock}
  \begin{block}{Direct computation}
    The probability of reaching $N$ is $\sum_n\Prob(\text{the game ends at round $n$ with fortune $N$})$, a sum over paths of every length $n$ that stay strictly between $0$ and $N$ until the last step. The summands have no closed form and there are infinitely many of them; the event is not determined by the fortune at any fixed time.
  \end{block}
  \words{The question concerns where the path leaves an interval rather than its value at a fixed time; the same question arises for an insurer's reserves and for the length of a queue.}
\end{frame}

% ---------------------------------------------------------------------------
\begin{frame}{Risk of Ruin: The Setting}
  \label{fr:ruinsetting}
  \small
  \begin{block}{Definition}
    With $\xi_1,\xi_2,\dots$ the Bernoulli process of slide~\ref{fr:walk}, the fortune is the walk started at $i$ and stopped on exit from $\{1,\dots,N-1\}$:
    \[
      X_n=i+S_n\ \text{ for } n\le T,\qquad T=\min\{n\ge0 : X_n\in\{0,N\}\},\qquad
      h_i=\Prob_i(X_T=N).
    \]
    The function $h\colon\{0,1,\dots,N\}\to[0,1]$ is the \textbf{probability of reaching the target} from fortune $i$; $1-h_i$ is the \textbf{probability of ruin}.
  \end{block}
  \words{The fortune is a random walk on $\{0,\dots,N\}$ stopped at the endpoints, and $h_i$ is the probability that the walk started at $i$ reaches $N$ before $0$.}
  \begin{center}
    \begin{tikzpicture}[x=1.35cm,y=0.85cm,font=\small,>=Stealth]
      \foreach \k/\l in {0/0,1/1,2/2,3/{i-1},4/i,5/{i+1},6/{N-1},7/N}{
        \node[circle,draw,minimum size=7mm,inner sep=0pt] (s\k) at (\k,0) {$\l$};}
      \node[circle,draw,minimum size=7mm,inner sep=0pt,fill=amber!20] at (0,0) {$0$};
      \node[circle,draw,minimum size=7mm,inner sep=0pt,fill=leaf!20] at (7,0) {$N$};
      \foreach \a/\b in {1/2,4/5,5/6}{
        \draw[->,accent] (s\a) to[bend left=35] node[above,font=\scriptsize] {$p$} (s\b);
        \draw[->,ink] (s\b) to[bend left=35] node[below,font=\scriptsize] {$q$} (s\a);}
      \draw[->,ink] (s1) to[bend left=35] node[below,font=\scriptsize] {$q$} (s0);
      \draw[->,accent] (s6) to[bend left=35] node[above,font=\scriptsize] {$p$} (s7);
      \draw[->,accent] (s3) to[bend left=35] node[above,font=\scriptsize] {$p$} (s4);
      \draw[->,ink] (s4) to[bend left=35] node[below,font=\scriptsize] {$q$} (s3);
      \node at (2.5,0) {$\cdots$}; \node at (5.5,0) {$\cdots$};
      \node[below=6mm,amber!80!black] at (s0) {ruin};
      \node[below=6mm,leaf!80!black] at (s7) {target};
    \end{tikzpicture}
  \end{center}
  \begin{itemize}
    \item $T$ is a random variable $\Omega\to\{0,1,2,\dots\}\cup\{\infty\}$; that $T<\infty$ with probability one follows from the analysis below (the expected duration is finite).
  \end{itemize}
\end{frame}

% ---------------------------------------------------------------------------
\begin{frame}{First-Step Analysis: The Ruin Equation}
  \label{fr:ruineq}
  \small
  \begin{block}{Theorem}
    The function $h$ satisfies
    \[
      h_i=p\,h_{i+1}+q\,h_{i-1}\quad(0<i<N),\qquad h_0=0,\qquad h_N=1 .
    \]
  \end{block}
  \words{After one round the fortune is $i+1$ or $i-1$ and the game continues as if started from there, so $h_i$ is a weighted average of $h_{i+1}$ and $h_{i-1}$.}
  \begin{block}{Proof}
    $h_0=0$ and $h_N=1$ because in these states the game has already ended. For $0<i<N$, condition on the first round (law of total probability with the two rows $\{\xi_1=+1\}$ and $\{\xi_1=-1\}$):
    \[
      h_i=\Prob_i(X_T=N\mid\xi_1=+1)\,p+\Prob_i(X_T=N\mid\xi_1=-1)\,q .
    \]
    On the row $\{\xi_1=+1\}$ the fortune is $i+1$ after round one and, by independent increments, the remaining rounds $\xi_2,\xi_3,\dots$ form a Bernoulli process independent of $\xi_1$: the conditional probability of reaching $N$ is $h_{i+1}$. Likewise $h_{i-1}$ on the other row. \hfill$\square$
  \end{block}
  \begin{itemize}
    \item This is the argument of the Conditional Expectation lecture (Application 2) and of slide~\ref{fr:firststep}, applied to a probability; it gives $N-1$ linear equations in $h_1,\dots,h_{N-1}$.
  \end{itemize}
\end{frame}

% ---------------------------------------------------------------------------
\begin{frame}{The Fair Game: $h_i=i/N$}
  \label{fr:ruinfair}
  \small
  \begin{block}{Theorem ($p=q=\tfrac12$)}
    $h_i=\dfrac iN$ for $i=0,\dots,N$. \quad Probability of ruin: $1-\dfrac iN$.
  \end{block}
  \words{In a fair game the probability of reaching $N$ from $i$ equals the fraction $i/N$ of the target already held.}
  \begin{block}{Proof}
    The equation $h_i=\tfrac12h_{i+1}+\tfrac12h_{i-1}$ says $h_{i+1}-h_i=h_i-h_{i-1}$: consecutive differences are equal, so $h$ is linear in $i$, $h_i=A+Bi$. Then $h_0=0$ gives $A=0$ and $h_N=1$ gives $B=\tfrac1N$.
    Uniqueness: if $h,h'$ both solve the system, $d=h-h'$ is linear with $d_0=d_N=0$, hence $d=0$. \hfill$\square$
  \end{block}
  \begin{block}{Second proof: the fair game preserves expected fortune}
    Each round has mean gain $0$, so $\E[X_T]=X_0=i$ (this needs an argument for the random time $T$, which we omit). But $X_T\in\{0,N\}$, so $\E[X_T]=N\,h_i$. Hence $h_i=i/N$.
  \end{block}
  % python3 fractions: win(3,10,1/2)=3/10; win(50,100,1/2)=1/2; win(1,10,1/2)=1/10
  \begin{itemize}
    \item $\pounds50$ to $\pounds100$ at a fair game: $\tfrac12$. $\pounds3$ to $\pounds10$: $\tfrac3{10}$. $\pounds1$ to $\pounds10$: $\tfrac1{10}$, and ruin $\tfrac9{10}$.
  \end{itemize}
\end{frame}

% ---------------------------------------------------------------------------
\begin{frame}{The Unfair Game: $h_i=\dfrac{r^i-1}{r^N-1}$ with $r=q/p$}
  \label{fr:ruinunfair}
  \small
  \begin{block}{Theorem ($p\ne q$)}
    With $r=q/p\ne1$: \quad $h_i=\dfrac{r^i-1}{r^N-1}$, \quad probability of ruin $1-h_i=\dfrac{r^N-r^i}{r^N-1}$.
  \end{block}
  \words{The probabilities now depend on $i$ through the geometric sequence $r^i$ rather than linearly, so a small disadvantage per round accumulates over the many rounds needed to reach the target.}
  \begin{block}{Proof}
    Try $h_i=r^i$ in $h_i=p\,h_{i+1}+q\,h_{i-1}$: the right side is $r^{i-1}(pr^2+q)$, and $pr^2+q=\tfrac{q^2}{p}+q=\tfrac{q(q+p)}{p}=r$, so it equals $r^i$. Constants also solve the equation. Hence $h_i=A+Br^i$ solves it for every $A,B$, and since a solution is fixed by $h_0,h_1$ (the equation determines $h_{i+1}$ from $h_i,h_{i-1}$), these are all the solutions. Now $h_0=0$ gives $A=-B$ and $h_N=1$ gives $B(r^N-1)=1$. \hfill$\square$
  \end{block}
  % python3 fractions (verified against sympy solve of the N=6 linear system for p=1/2 and p=2/5):
  %   win(3,10,2/5)=2432/58025=0.0419; win(5,10,2/5)=32/275=0.1164; win(5,10,3/5)=243/275=0.8836
  \begin{exampleblock}{$p=0.4$, $r=1.5$}
    From $\pounds5$ to $\pounds10$: $h_5=\dfrac{1.5^5-1}{1.5^{10}-1}=\dfrac{32}{275}\approx0.116$, against $\tfrac12$ in the fair game. From $\pounds3$ to $\pounds10$: $h_3\approx0.042$, against $0.3$.
  \end{exampleblock}
\end{frame}

% ---------------------------------------------------------------------------
\begin{frame}{Expected Duration of the Game}
  \label{fr:duration}
  \small
  Let $D_i=\E_i[T]$, so $D\colon\{0,\dots,N\}\to[0,\infty)$.
  \begin{block}{Theorem}
    $D_i=1+p\,D_{i+1}+q\,D_{i-1}$ for $0<i<N$, $D_0=D_N=0$, and the solution is
    \[
      D_i=i\,(N-i)\quad(p=q),\qquad
      D_i=\frac{i}{q-p}-\frac{N}{q-p}\cdot\frac{1-r^i}{1-r^N}\quad(p\ne q,\ r=q/p).
    \]
  \end{block}
  \words{In the fair game started from the midpoint the expected duration is $N^2/4$ rounds, the square of the distance to the boundary, as on slide~\ref{fr:exitsol}.}
  \begin{block}{Proof}
    The equation is first-step analysis with one round paid before moving to $i\pm1$, as on slide~\ref{fr:firststep}. For $p=q$ it is the equation of slide~\ref{fr:exitsol} on $\{0,\dots,N\}$; $D_i=i(N-i)$ satisfies it since $\tfrac12\big((i+1)(N-i-1)+(i-1)(N-i+1)\big)=i(N-i)-1$, and it vanishes at $0$ and $N$.
    For $p\ne q$: $D_i=\tfrac{i}{q-p}$ is one solution of the equation with the $+1$ ($\tfrac{p(i+1)+q(i-1)}{q-p}=\tfrac{i-(q-p)}{q-p}=\tfrac i{q-p}-1$); add the general solution $A+Br^i$ of the homogeneous equation and fit $D_0=D_N=0$. Uniqueness as before. \hfill$\square$
  \end{block}
  % python3: dur(50,100,1/2)=2500; dur(50,100,18/38)=940.26; dur(3,10,1/2)=21; dur(5,10,2/5)=211/11=19.18
  \begin{itemize}
    \item Fair, $\pounds50$ to $\pounds100$: $2500$ rounds on average. Roulette ($p=\tfrac{18}{38}$), same stakes: about $940$ rounds. The unfavourable game is both more likely to end in ruin and shorter on average.
  \end{itemize}
\end{frame}


% ===========================================================================
\begin{frame}{The Casino's Edge in Numbers}
  \label{fr:casino}
  \small
  \begin{center}
      % python3 fractions, win() and dur() above:
      % (50,100,1/2): 0.5, 2500   (50,100,49/100): 0.1192, 1904   (50,100,18/38): 0.00513, 940
      % (10,20,18/38): 0.2585, 91.8  (5,10,18/38): 0.3713, 24.5   (1,2,18/38): 0.4737, 1
      \renewcommand{\arraystretch}{1.25}
      \begin{tabular}{@{}lccc@{}}
        \toprule
        $i\to N$ (in units of one stake) & $p$ & $\Prob(\text{target})$ & $\E[T]$ \\
        \midrule
        $50\to100$ & $0.5$ & $0.500$ & $2500$ \\
        $50\to100$ & $0.49$ & $0.119$ & $1904$ \\
        $50\to100$ & $18/38$ & $0.005$ & $940$ \\
        $10\to20$ & $18/38$ & $0.259$ & $92$ \\
        $5\to10$ & $18/38$ & $0.371$ & $24$ \\
        $1\to2$ & $18/38$ & $0.474$ & $1$ \\
        \bottomrule
      \end{tabular}
  \end{center}
  \words{At a $1\%$ disadvantage, doubling $\pounds50$ one pound at a time succeeds with probability about $0.12$; the last three rows are $\pounds50\to\pounds100$ with stakes of $\pounds5$, $\pounds10$ and $\pounds50$, and with fewer rounds the disadvantage accumulates less.}
\end{frame}

% ---------------------------------------------------------------------------
\begin{frame}{Why Doubling Fails}
  \label{fr:doubling}
  \small
      \begin{block}{Doubling (the ``martingale'' system)}
        Bet $1$; after each loss double the stake; after a win collect the net gain of $\pounds1$ and restart. With capital $2^k-1$ the gambler can absorb $k-1$ consecutive losses; $k$ losses in a row exhaust the capital.
        Per cycle: gain $+1$ with probability $1-q^k$, lose $2^k-1$ with probability $q^k$:
        \[
          \E[\text{gain}]=(1-q^k)-q^k(2^k-1)=1-(2q)^k .
        \]
        This is $0$ if $q=\tfrac12$ and negative if $q>\tfrac12$. The strategy changes the distribution of the outcome, to many small wins and a rare large loss, but not its mean.
        % python3: q=20/38,k=7: 1-(2q)^7=-0.432; q^7=0.0112; loss 127; fair: 0
      \end{block}
      \begin{itemize}
        \item Roulette, $k=7$ (capital $\pounds127$): win $\pounds1$ with probability $0.989$, lose $\pounds127$ with probability $0.011$; mean $-\pounds0.43$ per cycle.
      \end{itemize}
  \words{Doubling produces many small wins and a rare large loss; the mean gain per cycle is negative whenever each bet has negative mean.}
\end{frame}

% ---------------------------------------------------------------------------
\begin{frame}{The Infinite Target: An Opponent With Unlimited Capital}
  \label{fr:infinite}
  \small
  Let $N\to\infty$ with $i$ fixed: the gambler plays on until ruin or forever.
  \begin{block}{Theorem}
    Let $\rho_i=\Prob_i(\text{ruin, ever})=\lim_{N\to\infty}(1-h_i)$. Then
    \[
      \rho_i=1\ \ (p\le q),\qquad \rho_i=\Big(\frac qp\Big)^{i}\ \ (p>q).
    \]
    In the fair case, moreover, $\E[\text{time to ruin}]=\lim_N i(N-i)=\infty$.
  \end{block}
  \words{Against an opponent with unlimited capital, a fair or unfavourable game ends in ruin with probability one; in a favourable game the probability of ruin is $(q/p)^i$, which decreases geometrically in the initial fortune $i$.}
  \begin{block}{Proof}
    Fair: $1-h_i=1-\tfrac iN\to1$. Unfair with $r=q/p$: $1-h_i=\dfrac{r^N-r^i}{r^N-1}$. If $r>1$ ($p<q$), divide by $r^N$: $\to\dfrac{1-0}{1-0}=1$. If $r<1$ ($p>q$), $r^N\to0$ and $1-h_i\to\dfrac{0-r^i}{0-1}=r^i$. \hfill$\square$
  \end{block}
\end{frame}

% ---------------------------------------------------------------------------
\begin{frame}{The Infinite Target: Three Regimes}
  \label{fr:regimes}
  \small
  % python3: (2/3)^5 = 0.1317; i/(q-p) for i=50, p=18/38: 950
  \begin{center}
    \footnotesize
    \renewcommand{\arraystretch}{1.15}
    \begin{tabular}{@{}lcc>{\raggedright\arraybackslash}p{5.8cm}@{}}
      \toprule
      game & $\Prob(\text{ruin})$ from $i$ & $\E[\text{time to ruin}]$ & the walk $i+S_n$ \\
      \midrule
      unfavourable, $p<q$ & $1$ & $\dfrac{i}{q-p}<\infty$ & $\E[\xi_n]<0$: reaches $0$ with probability $1$, in finite mean time \\
      fair, $p=q$ & $1$ & $\infty$ & $\E[\xi_n]=0$: reaches $0$ with probability $1$, but with infinite mean time \\
      favourable, $p>q$ & $(q/p)^i<1$ & $\infty$ & $\E[\xi_n]>0$: reaches $0$ with probability $(q/p)^i<1$ \\
      \bottomrule
    \end{tabular}
  \end{center}
  \begin{itemize}
    \item $p=0.6$, $\pounds5$ in hand: ruin with probability $(\tfrac23)^5\approx0.13$; otherwise the fortune drifts to infinity.
    \item Roulette, $\pounds50$ in hand, no target: ruin certain, after $\dfrac{50}{2/38}=950$ spins on average.
    \item Fair game: ruin is certain, but the expected time to ruin is infinite. This is the statement that the fair walk reaches any level with probability one but with infinite expected time; it reappears as \emph{null recurrence} on slide~\ref{fr:walkZ}.
    \item $\rho_i$ for $p>q$ is the probability that a walk with upward drift ever falls by $i$ below its start: the \textbf{hitting probability} of level $-i$. Hitting probabilities are the subject of slides~\ref{fr:hitting}--\ref{fr:minimal}, where the ruin problem returns as a special case.
  \end{itemize}
  \words{The three regimes differ in whether ruin is certain and, if it is, whether the expected time to ruin is finite.}
\end{frame}

% ---------------------------------------------------------------------------
\begin{frame}{Return to the Opening Problem}
  \label{fr:ruinreturn}
  \small
  % python3 fractions: win(50,100,18/38)=0.005127; dur=940.26; win(50,60,18/38)=0.3475, dur=553.8;
  %   win(50,55,18/38)=0.5892, dur=334.2; win(50,51,18/38)=0.8995, dur=78.4; one bet of 50: 18/38=0.4737
  Red at roulette, $p=\tfrac{18}{38}$, $r=\tfrac{20}{18}=\tfrac{10}9$. From $\pounds50$, one pound per spin:
  \begin{center}
    \renewcommand{\arraystretch}{1.25}
    \begin{tabular}{@{}lccc@{}}
      \toprule
      target $N$ & $h_{50}=\dfrac{r^{50}-1}{r^{N}-1}$ & ruin & $\E[T]$ (spins) \\
      \midrule
      $100$ (double) & $0.005$ & $0.995$ & $940$ \\
      $60$ (win $\pounds10$) & $0.348$ & $0.652$ & $554$ \\
      $55$ (win $\pounds5$) & $0.589$ & $0.411$ & $334$ \\
      $51$ (win $\pounds1$) & $0.900$ & $0.100$ & $78$ \\
      \bottomrule
    \end{tabular}
  \end{center}
  \begin{itemize}
    \item Doubling $\pounds50$ at $\pounds1$ a spin succeeds with probability about $1/195$, and the game lasts about $940$ spins on average. Staking all $\pounds50$ on a single spin succeeds with probability $\tfrac{18}{38}\approx0.474$.
    \item The plan ``win $\pounds1$ and stop'' fails with probability $0.1005$, and each failure costs $\pounds50$; the mean gain is $0.8995\cdot1-0.1005\cdot50\approx-\pounds4.12$, which equals $(p-q)\,\E[T]=-\tfrac2{38}\cdot78.4$. No staking plan has positive mean gain, since every plan is a sum of bets each of mean $-\tfrac2{38}$ per pound staked.
  \end{itemize}
  \words{The entries are the values of $h_{50}$ and $D_{50}$ given by the formulas of slides~\ref{fr:ruinunfair} and~\ref{fr:duration} for each target $N$.}
\end{frame}


% ===========================================================================
\begin{frame}{Markov Chains: The Problem}
  \label{fr:mcprob}
  \small
  A town's weather is $\st{Sunny}$, $\st{Cloudy}$ or $\st{Rainy}$ each day, and tomorrow's weather depends on today's according to the table below. It is sunny today.
  \begin{columns}[T]
    \begin{column}{0.56\textwidth}
      \begin{enumerate}
        \item What is the probability that it is rainy a week from now?
        \item In the long run, what fraction of days are sunny?
        \item How many days until the next rainy day, on average?
      \end{enumerate}
      \begin{block}{Direct computation}
        Question 1 is a sum over the $3^6=729$ possible weather sequences for the six intermediate days. Question 2 involves a limit over sequences of every length. Question 3 concerns a waiting time whose distribution has no simple closed form.
      \end{block}
      \words{In this model tomorrow's weather depends on the past only through today's; this property reduces all three questions to computations with the $3\times3$ matrix.}
    \end{column}
    \begin{column}{0.42\textwidth}
      \centering
      \renewcommand{\arraystretch}{1.3}
      \begin{tabular}{@{}l|ccc@{}}
        today $\backslash$ tomorrow & $\st S$ & $\st C$ & $\st R$ \\ \hline
        $\st S$ & $0.7$ & $0.2$ & $0.1$ \\
        $\st C$ & $0.3$ & $0.4$ & $0.3$ \\
        $\st R$ & $0.2$ & $0.4$ & $0.4$ \\
      \end{tabular}
      \par\medskip
      {\footnotesize Row $\st S$ reads: after a sunny day, tomorrow is sunny with probability $0.7$, cloudy $0.2$, rainy $0.1$. Each row sums to $1$.}
    \end{column}
  \end{columns}
\end{frame}

% ---------------------------------------------------------------------------
\begin{frame}{Markov Chain: Definition}
  \label{fr:mcdef}
  \small
  \begin{block}{Definition}
    Let $S$ be a finite or countable set (the \textbf{state space}). A process $(X_n)_{n\ge0}$ with $X_n\colon\Omega\to S$ is a \textbf{Markov chain} if for all $n\ge0$ and all states $i_0,\dots,i_{n-1},i,j\in S$ with $\Prob(X_0=i_0,\dots,X_{n-1}=i_{n-1},X_n=i)>0$,
    \[
      \Prob(X_{n+1}=j\mid X_0=i_0,\dots,X_{n-1}=i_{n-1},\ X_n=i)=\Prob(X_{n+1}=j\mid X_n=i)\qquad\text{(Markov property)},
    \]
    and it is \textbf{homogeneous} if the right side does not depend on $n$. Then the \textbf{transition probabilities} are the function
    \[
      p\colon S\times S\to[0,1],\qquad p(i,j)=\Prob(X_{n+1}=j\mid X_n=i).
    \]
  \end{block}
  \words{The conditional distribution of the next state given the entire past depends only on the present state, and in a homogeneous chain it is the same at every time.}
  \begin{itemize}
    \item The two arguments of $p$ play different roles: $i$ is the \emph{current} state (the conditioning event $\{X_n=i\}$, the row); $j$ is the \emph{next} state (the column). $p(i,j)$ is the probability of the step $i\to j$.
    \item Random walk: $S=\Z$, $p(i,i+1)=p$, $p(i,i-1)=q$, all other $p(i,j)=0$. The Markov property holds because $S_{n+1}=S_n+\xi_{n+1}$ with $\xi_{n+1}$ independent of the past. Gambler's ruin: the same on $S=\{0,\dots,N\}$ with $p(0,0)=p(N,N)=1$.
    \item All chains below are homogeneous.
  \end{itemize}
\end{frame}

% ---------------------------------------------------------------------------
\begin{frame}{The Transition Matrix: A Table Whose Rows Are Conditional Distributions}
  \label{fr:tm}
  \small
  \begin{columns}[T]
    \begin{column}{0.60\textwidth}
      \begin{block}{Definition}
        The \textbf{transition matrix} is $p$ written as a table, $\PP=\big(p(i,j)\big)_{i,j\in S}$: row $i$, column $j$. It is a \textbf{stochastic matrix}:
        \[
          p(i,j)\ge0\quad\text{and}\quad\sum_{j\in S}p(i,j)=1\ \text{ for every } i\in S .
        \]
        The \textbf{initial distribution} is the row vector $\lambda\colon S\to[0,1]$, $\lambda_i=\Prob(X_0=i)$, $\sum_i\lambda_i=1$.
      \end{block}
      \words{Row $i$ of $\PP$ is the conditional pmf of $X_{n+1}$ given $X_n=i$, as a row of the two-dice table is.}
      \begin{itemize}
        \item Row $i$, as a function $j\mapsto p(i,j)$, is $p_{X_{n+1}\mid X_n}(\,\cdot\mid i)$: non-negative, sums to $1$. Column sums are unconstrained.
        \item $\E[f(X_{n+1})\mid X_n=i]=\sum_j p(i,j)f(j)$: the row average of $f$ (compare $\E[X\mid Y=y]$, lecture 5).
      \end{itemize}
    \end{column}
    \begin{column}{0.38\textwidth}
      \centering
      \begin{tikzpicture}[x=1.0cm,y=0.7cm,font=\small]
        \fill[shade] (-0.5,-1.5) rectangle (2.5,-2.5);
        \foreach \j/\l in {0/{$\st S$},1/{$\st C$},2/{$\st R$}}{\node[soft] at (\j,0.5) {\l};}
        \foreach \i/\l in {1/{$\st S$},2/{$\st C$},3/{$\st R$}}{\node[soft] at (-1,-\i) {\l};}
        \node at (0,-1) {0.7}; \node at (1,-1) {0.2}; \node at (2,-1) {0.1};
        \node at (0,-2) {0.3}; \node at (1,-2) {0.4}; \node at (2,-2) {0.3};
        \node at (0,-3) {0.2}; \node at (1,-3) {0.4}; \node at (2,-3) {0.4};
        \draw[soft] (-0.5,-0.5) grid[xstep=1,ystep=1] (2.5,-3.5);
        \node[soft,font=\scriptsize] at (1,1.1) {next state $j$};
        \node[soft,font=\scriptsize,rotate=90] at (-1.7,-2) {current state $i$};
      \end{tikzpicture}
      \par\medskip
      {\footnotesize The weather chain. The shaded row is $p_{X_{n+1}\mid X_n}(\,\cdot\mid\st C)$: the distribution of tomorrow given a cloudy today.}
    \end{column}
  \end{columns}
\end{frame}

% ---------------------------------------------------------------------------
\begin{frame}{Paths Have Product Probabilities}
  \label{fr:paths}
  \small
  \begin{block}{Theorem}
    For a Markov chain with initial distribution $\lambda$ and transition matrix $\PP$, and any states $i_0,i_1,\dots,i_n$,
    \[
      \Prob(X_0=i_0,\ X_1=i_1,\dots,X_n=i_n)=\lambda_{i_0}\,p(i_0,i_1)\,p(i_1,i_2)\cdots p(i_{n-1},i_n).
    \]
    Conversely, a process whose joint pmfs have this form is a Markov chain with transition matrix $\PP$.
  \end{block}
  \words{The probability of a row of the table of paths is the start probability times the product of the step probabilities along the row; $\lambda$ and $\PP$ together determine every finite-dimensional distribution.}
  \begin{block}{Proof}
    Multiplication rule: $\Prob(X_0=i_0,\dots,X_n=i_n)=\Prob(X_0=i_0)\prod_{k=1}^{n}\Prob(X_k=i_k\mid X_0=i_0,\dots,X_{k-1}=i_{k-1})$, and each factor is $p(i_{k-1},i_k)$ by the Markov property. For the converse, divide two such products to get the conditional probability. \hfill$\square$
  \end{block}
  % python3: 0.2*0.3 = 0.06; 0.1*0.4*0.2 = 0.008
  \begin{exampleblock}{Weather chain (slide~\ref{fr:tm}), sunny today}
    $\Prob(X_1=\st C,\ X_2=\st R\mid X_0=\st S)=p(\st S,\st C)\,p(\st C,\st R)=0.2\cdot0.3=0.06$; \quad $\Prob(X_1=\st R,X_2=\st R,X_3=\st S\mid X_0=\st S)=0.1\cdot0.4\cdot0.2=0.008$.
  \end{exampleblock}
\end{frame}

% ---------------------------------------------------------------------------
\begin{frame}{What the Markov Property Is Not}
  \label{fr:mcnot}
  \small
  \begin{itemize}
    \item \textbf{It is not independence.} $X_{n+1}$ may depend strongly on $X_n$; the property says that the earlier states $X_0,\dots,X_{n-1}$ carry no further information once $X_n$ is known.
    \item \textbf{It does not say the past is irrelevant.} In general $\Prob(X_{n+1}=j\mid X_{n-1}=k)\ne\Prob(X_{n+1}=j)$; the past is irrelevant only conditionally on the present.
    \item \textbf{Not every process has it.} Toss a coin and let $X_n$ be the number of heads among tosses $n-1$ and $n$. Given $X_n=1$, the distribution of $X_{n+1}$ depends on which of the two tosses was the head: $\Prob(X_{n+1}=2\mid X_n=1,X_{n-1}=0)=\tfrac12$ but $\Prob(X_{n+1}=2\mid X_n=1,X_{n-1}=2)=0$. It is not Markov; the process recording the last two tosses is.
    \item \textbf{It depends on what is recorded as the state.} A process that is not Markov often becomes Markov when the state is enlarged to include the relevant part of the past (the last two tosses; position and velocity; today's weather and yesterday's).
  \end{itemize}
  \begin{block}{Test for the Markov property}
    Check that $\Prob(X_{n+1}=j\mid X_n=i,\ \text{anything about } X_0,\dots,X_{n-1})$ does not change with the ``anything''. In models of the form $X_{n+1}=F(X_n,\xi_{n+1})$ with $\xi_{n+1}$ independent of $X_0,\dots,X_n$ it holds automatically; the random walk, the Ehrenfest urn and the walks on graphs below are all of this form.
  \end{block}
  \words{The Markov property states that, for predicting the future, the present state is a sufficient summary of the past.}
\end{frame}

% ---------------------------------------------------------------------------
\begin{frame}{A Gallery of Chains (1)}
  \label{fr:gallery1}
  \small
  \begin{columns}[T]
    \begin{column}{0.48\textwidth}
      \begin{block}{Vowels and consonants}
        Model a text as a sequence of $\st V$ and $\st C$. After a vowel, the next letter is a vowel with probability $0.2$; after a consonant, with probability $0.5$.
        \[
          S=\{\st V,\st C\},\qquad \PP=\begin{pmatrix}0.2&0.8\\0.5&0.5\end{pmatrix}.
        \]
        This two-state chain is used in the Entropy lecture for the entropy-rate example.
      \end{block}
      \begin{block}{Gambler's ruin}
        $S=\{0,\dots,N\}$, $p(i,i+1)=p$, $p(i,i-1)=q$ for $0<i<N$, and $p(0,0)=p(N,N)=1$: the diagram of slide~\ref{fr:ruinsetting}. States $0$ and $N$ are \textbf{absorbing}: once the chain enters one of them it remains there.
      \end{block}
    \end{column}
    \begin{column}{0.48\textwidth}
      \begin{block}{The Ehrenfest urn (gas between two chambers)}
        $N$ numbered balls in two urns. Each step, pick a ball uniformly at random and move it to the other urn. $X_n$ is the number of balls in urn~$1$: $S=\{0,\dots,N\}$,
        \[
          p(i,i-1)=\frac iN,\qquad p(i,i+1)=\frac{N-i}{N}.
        \]
        The more balls an urn holds, the more likely the next move is out of it.
        % python3 sympy: N=4 matrix rows: (0,1,0,0,0),(1/4,0,3/4,0,0),(0,1/2,0,1/2,0),(0,0,3/4,0,1/4),(0,0,0,1,0)
        \[
          N=4:\quad \PP=\begin{pmatrix}0&1&0&0&0\\ \tfrac14&0&\tfrac34&0&0\\ 0&\tfrac12&0&\tfrac12&0\\ 0&0&\tfrac34&0&\tfrac14\\ 0&0&0&1&0\end{pmatrix}.
        \]
      \end{block}
    \end{column}
  \end{columns}
\end{frame}

% ---------------------------------------------------------------------------
\begin{frame}{A Gallery of Chains (2): Walks on Graphs}
  \label{fr:gallery2}
  \small
  \begin{columns}[T]
    \begin{column}{0.55\textwidth}
      \begin{block}{Random walk on a graph}
        Vertices $S$; from vertex $i$ move to one of its $d(i)$ neighbours uniformly: $p(i,j)=1/d(i)$ if $j$ is a neighbour of $i$, else $0$.
      \end{block}
      \begin{itemize}
        \item \textbf{Cycle} of length $L$: $S=\{0,\dots,L-1\}$, $p(i,i\pm1\bmod L)=\tfrac12$. This is the random walk on $\{0,\dots,L-1\}$ with the two ends joined.
        \item \textbf{Cube}: $S=\{0,1\}^3$, the $8$ corners; neighbours differ in one coordinate, so $p(i,j)=\tfrac13$ for each of the $3$ neighbours. A move flips one coordinate.
        \item \textbf{The web}: pages as vertices, links as directed edges, $p(i,j)=1/(\text{number of links out of } i)$ for each page $j$ that $i$ links to. This is the \textbf{random surfer} of slides~\ref{fr:pagerank}--\ref{fr:pagerank2}.
      \end{itemize}
      \words{On a graph, the transition matrix is the adjacency table with each row rescaled to sum to $1$.}
    \end{column}
    \begin{column}{0.43\textwidth}
      \centering
      \begin{tikzpicture}[scale=1.05,font=\scriptsize]
        \coordinate (000) at (0,0); \coordinate (100) at (1.4,0); \coordinate (010) at (0,1.4); \coordinate (110) at (1.4,1.4);
        \coordinate (001) at (0.6,0.5); \coordinate (101) at (2.0,0.5); \coordinate (011) at (0.6,1.9); \coordinate (111) at (2.0,1.9);
        \draw[soft] (000)--(100)--(110)--(010)--cycle; \draw[soft] (001)--(101)--(111)--(011)--cycle;
        \draw[soft] (000)--(001) (100)--(101) (110)--(111) (010)--(011);
        \foreach \v in {000,100,010,110,001,101,011,111}{\fill[ink] (\v) circle (1.6pt);}
        \fill[accent] (000) circle (2.6pt); \fill[amber] (111) circle (2.6pt);
        \node[below left] at (000) {$000$}; \node[above right] at (111) {$111$};
        \node[below right] at (100) {$100$}; \node[above left] at (010) {$010$}; \node[right] at (001) {$001$};
      \end{tikzpicture}
      \hspace{4mm}
      \begin{tikzpicture}[scale=0.9,font=\scriptsize]
        \foreach \k in {0,...,4}{\coordinate (c\k) at ({90+72*\k}:1);}
        \draw[soft] (c0)--(c1)--(c2)--(c3)--(c4)--cycle;
        \foreach \k in {0,...,4}{\fill[ink] (c\k) circle (1.6pt); \node at ({90+72*\k}:1.3) {$\k$};}
      \end{tikzpicture}
      \par\medskip
      {\footnotesize Left: the cube; from $000$, the three neighbours $100,010,001$. Right: the cycle $L=5$.}
    \end{column}
  \end{columns}
\end{frame}


% ===========================================================================
\begin{frame}{$n$-Step Transition Probabilities}
  \label{fr:nstep}
  \small
  As on slide~\ref{fr:firststep}, $\Prob_i(\,\cdot\,)=\Prob(\,\cdot\mid X_0=i)$ and $\E_i[\,\cdot\,]=\E[\,\cdot\mid X_0=i]$ refer to the chain started at $i$ (the case $\lambda=$ point mass at $i$).
  \begin{block}{Definition}
    The \textbf{$n$-step transition probabilities} are the function $p^{(n)}\colon S\times S\to[0,1]$,
    \[
      p^{(n)}(i,j)=\Prob_i(X_n=j)=\Prob(X_{m+n}=j\mid X_m=i)\quad\text{(any $m$, by homogeneity)},
    \]
    with $p^{(1)}=p$ and $p^{(0)}(i,j)=\ind\{i=j\}$. As before, $i$ is the current state (row) and $j$ the state $n$ steps later (column).
  \end{block}
  \words{The $n$-step probabilities form a table in which row $i$ is the conditional distribution of $X_{m+n}$ given $X_m=i$.}
  % python3: 0.7*0.1 + 0.2*0.3 + 0.1*0.4 = 0.17
  \begin{exampleblock}{Weather: $p^{(2)}(\st S,\st R)$ by hand}
    Sum over tomorrow's weather $k$: $p^{(2)}(\st S,\st R)=\sum_k p(\st S,k)\,p(k,\st R)=0.7\cdot0.1+0.2\cdot0.3+0.1\cdot0.4=0.17$. The three terms are the sequences $\st{SSR},\st{SCR},\st{SRR}$, each a product as on slide~\ref{fr:paths}.
  \end{exampleblock}
  \begin{itemize}
    \item This sum is the $(\st S,\st R)$ entry of $\PP\PP$; the next slide shows that this holds in general.
  \end{itemize}
\end{frame}

% ---------------------------------------------------------------------------
\begin{frame}{Chapman--Kolmogorov: $n$ Steps Is the $n$-th Power}
  \label{fr:ck}
  \small
  \begin{block}{Theorem (Chapman--Kolmogorov)}
    For all $m,n\ge0$ and $i,j\in S$: \quad $\displaystyle p^{(m+n)}(i,j)=\sum_{k\in S}p^{(m)}(i,k)\,p^{(n)}(k,j)$. \quad Equivalently, writing $\PP^{(n)}=\big(p^{(n)}(i,j)\big)_{i,j}$ for the table of $n$-step probabilities: $\PP^{(n)}=\PP^n$, the $n$-th matrix power.
  \end{block}
  \words{A path from $i$ to $j$ in $m+n$ steps is in some state $k$ at time $m$; summing over $k$ gives the formula, which has the form of a matrix product.}
  \begin{block}{Proof}
    Total probability with the rows $\{X_m=k\}$, $k\in S$:
    \[
      \Prob_i(X_{m+n}=j)=\sum_k\Prob_i(X_m=k)\,\Prob_i(X_{m+n}=j\mid X_m=k)=\sum_k p^{(m)}(i,k)\,p^{(n)}(k,j),
    \]
    the second factor by the Markov property and homogeneity (the past before time $m$, including $X_0=i$, drops out). With $m=1$ this is $\PP^{(n+1)}=\PP\,\PP^{(n)}$, so by induction $\PP^{(n)}=\PP^n$. \hfill$\square$
  \end{block}
  \begin{itemize}
    \item The case $m=1$ is first-step analysis and $n=1$ a last-step decomposition; numerically, $\PP^n$ is computed by repeated squaring. Rows of $\PP^n$ are conditional distributions, so $\PP^n$ is stochastic.
  \end{itemize}
\end{frame}

% ---------------------------------------------------------------------------
\begin{frame}{The Distribution at Time $n$: Row Vector Times $\PP^n$}
  \label{fr:lawn}
  \small
  \begin{block}{Theorem}
    If $X_0$ has distribution $\lambda$ (a row vector indexed by $S$), then $X_n$ has distribution $\lambda\PP^n$:
    \[
      \Prob(X_n=j)=\sum_{i\in S}\lambda_i\,p^{(n)}(i,j)=(\lambda\PP^n)_j .
    \]
  \end{block}
  \words{The distribution at time $n$ is the initial distribution multiplied by $\PP$ once for each step; each multiplication forms a weighted average of the rows of $\PP$ with the current probabilities as weights.}
  \begin{block}{Proof}
    Total probability over the rows $\{X_0=i\}$: $\Prob(X_n=j)=\sum_i\Prob(X_0=i)\Prob(X_n=j\mid X_0=i)$. \hfill$\square$
  \end{block}
  \begin{itemize}
    \item Every finite-dimensional distribution follows: $\Prob(X_m=i,\ X_{m+n}=j)=(\lambda\PP^m)_i\,p^{(n)}(i,j)$, and so on along any sequence of times.
    \item $\lambda\PP^n$ is a row vector, a distribution on $S$, while $\PP^n$ is a matrix, a table of conditional distributions. Multiplication by $\PP$ advances a distribution by one step in time.
  \end{itemize}
\end{frame}

% ---------------------------------------------------------------------------
\begin{frame}{Return to Question 1: Rainy a Week From a Sunny Day}
  \label{fr:q1}
  \small
  % python3 numpy: W=[[.7,.2,.1],[.3,.4,.3],[.2,.4,.4]]; matrix_power(W,2)[0]=(0.57,0.26,0.17);
  %   matrix_power(W,3)=[[.511,.286,.203],[.429,.322,.249],[.406,.332,.262]]
  %   matrix_power(W,7)[0]=(0.4637,0.3068,0.2296); matrix_power(W,20) rows all (0.4615,0.3077,0.2308)
  $\lambda=(1,0,0)$, so the distribution of $X_7$ is row $\st S$ of $\PP^7$.
  {\setlength{\arraycolsep}{3pt}
  \[
    \PP^2=\begin{pmatrix}0.57&0.26&0.17\\0.39&0.34&0.27\\0.34&0.36&0.30\end{pmatrix},\quad
    \PP^3=\begin{pmatrix}0.511&0.286&0.203\\0.429&0.322&0.249\\0.406&0.332&0.262\end{pmatrix},\quad
    \PP^7=\begin{pmatrix}0.4637&0.3068&0.2296\\0.4601&0.3083&0.2316\\0.4591&0.3087&0.2321\end{pmatrix}.
  \]}
  \begin{exampleblock}{Answer}
    $\Prob_{\st S}(X_7=\st R)=p^{(7)}(\st S,\st R)\approx0.230$. The sum over $729$ six-day sequences has been organised into six matrix multiplications.
  \end{exampleblock}
  \begin{itemize}
    \item The three rows of $\PP^7$ are almost equal: after a week the distribution of the weather depends very little on today's weather. The common limiting row, $(0.4615,0.3077,0.2308)$ to four places from $\PP^{20}$, is the subject of the next section.
    \item Comparing $\PP^2$ with $\PP$: after a rainy day the probability of sun is $0.2$ for the next day and $0.34$ for the day after; the rows of $\PP^n$ approach one another as $n$ grows.
  \end{itemize}
  \words{The distribution of the chain at any future time is obtained from the powers of $\PP$.}
\end{frame}

% ---------------------------------------------------------------------------
\begin{frame}{Powers of a $2\times2$ Chain: The Rows Converge}
  \label{fr:twobytwo}
  \small
  % python3 numpy: V=[[.2,.8],[.5,.5]]: V^2=[[.44,.56],[.35,.65]], V^3=[[.368,.632],[.395,.605]],
  %   V^10=[[.38462,.61538],[.38461,.61539]]; eigenvalues 1 and -0.3; 5/13=0.384615, 8/13=0.615385
  \begin{exampleblock}{Vowel/consonant chain}
    \[
      \PP=\begin{pmatrix}0.2&0.8\\0.5&0.5\end{pmatrix},\quad
      \PP^2=\begin{pmatrix}0.44&0.56\\0.35&0.65\end{pmatrix},\quad
      \PP^3=\begin{pmatrix}0.368&0.632\\0.395&0.605\end{pmatrix},\quad
      \PP^{10}=\begin{pmatrix}0.3846&0.6154\\0.3846&0.6154\end{pmatrix}.
    \]
    Both rows approach $(\tfrac5{13},\tfrac8{13})=(0.384615\ldots,\,0.615385\ldots)$; after ten letters the probability of a vowel no longer depends appreciably on whether the first letter was a vowel or a consonant.
  \end{exampleblock}
  \begin{block}{Exact formula}
    $\PP$ has eigenvalues $1$ and $0.2+0.5-1=-0.3$ (trace minus $1$), and
    \[
      p^{(n)}(\st V,\st V)=\tfrac5{13}+\tfrac8{13}(-0.3)^n,\qquad p^{(n)}(\st C,\st V)=\tfrac5{13}-\tfrac5{13}(-0.3)^n .
    \]
    For $n=1$ these give $\tfrac5{13}-\tfrac{2.4}{13}=0.2$ and $\tfrac5{13}+\tfrac{1.5}{13}=0.5$, as required. The difference from the limit is multiplied by $-0.3$ at each step.
  \end{block}
  \words{The remaining sections identify the limit row, give conditions under which it exists, and estimate the speed of convergence.}
\end{frame}


% ===========================================================================
\begin{frame}{Stationary Distribution: Definition}
  \label{fr:statdef}
  \small
  \begin{block}{Definition}
    A \textbf{stationary distribution} of a chain with transition matrix $\PP$ on $S$ is a function $\pi\colon S\to[0,1]$ with $\sum_{i\in S}\pi_i=1$ and
    \[
      \pi\PP=\pi,\qquad\text{i.e.}\qquad \sum_{i\in S}\pi_i\,p(i,j)=\pi_j\ \text{ for every } j\in S .
    \]
  \end{block}
  \words{A stationary distribution is one that the chain preserves: if $X_0$ has distribution $\pi$ then every $X_n$ has distribution $\pi$.}
  \begin{block}{Theorem}
    If $X_0\sim\pi$ and $\pi$ is stationary, then $X_n\sim\pi$ for every $n$, and $(X_n,X_{n+1})$ has the same joint distribution for every $n$.
  \end{block}
  \begin{block}{Proof}
    $\pi\PP^n=(\pi\PP)\PP^{n-1}=\pi\PP^{n-1}=\dots=\pi$, then the previous slide. $\Prob(X_n=i,X_{n+1}=j)=\pi_i\,p(i,j)$ does not involve $n$. \hfill$\square$
  \end{block}
\end{frame}

% ---------------------------------------------------------------------------
\begin{frame}{Stationary Distribution: Two Readings and Existence}
  \label{fr:statexist}
  \small
  \begin{block}{Balance}
    $\sum_i\pi_i\,p(i,j)$ is the probability mass flowing into $j$ in one step from the distribution $\pi$; stationarity says it equals the mass $\pi_j$ already there, for every $j$. Mass leaves $j$ at rate $\pi_j\sum_{k}p(j,k)=\pi_j$, so inflow equals outflow at every state.
  \end{block}
  \begin{block}{Left eigenvector}
    $\pi\PP=\pi$ says $\pi$ is a left eigenvector of $\PP$ for the eigenvalue $1$, normalised to sum to $1$. Every stochastic matrix has the eigenvalue $1$, because the column vector $\mathbf 1$ of ones satisfies $\PP\mathbf 1=\mathbf 1$ (row sums), and left and right eigenvalues coincide.
  \end{block}
  \begin{block}{Theorem (existence)}
    Every Markov chain on a finite state space has at least one stationary distribution.
  \end{block}
  \words{A left eigenvector for the eigenvalue $1$ exists by the argument above; that it can be chosen with non-negative entries is the content of the theorem.}
  \begin{itemize}
    \item The theorem is not proved here (Perron--Frobenius, or a fixed-point argument). Uniqueness can fail, as for the gambler's ruin chain, which has infinitely many. On an infinite $S$ existence can fail: the random walk on $\Z$ has none, since any solution of $\pi\PP=\pi$ is constant and cannot sum to $1$.
  \end{itemize}
\end{frame}

% ---------------------------------------------------------------------------
\begin{frame}{Computing $\pi$: A Linear System}
  \label{fr:statcompute}
  \small
  $\pi\PP=\pi$ is $|S|$ linear equations in $|S|$ unknowns, but they are dependent (they sum to the identity $\sum_j\pi_j=\sum_i\pi_i$); drop one and add $\sum_i\pi_i=1$.
  % python3 sympy: solve(pi W = pi, sum=1) -> (6/13, 4/13, 3/13) = (0.4615, 0.3077, 0.2308); matches W^20 rows
  \begin{exampleblock}{Weather, question 2: the long-run fraction of sunny days}
    Unknowns $\pi_{\st S},\pi_{\st C},\pi_{\st R}$. Columns $\st S$ and $\st C$ of $\pi\PP=\pi$, and normalisation:
    \[
      0.7\pi_{\st S}+0.3\pi_{\st C}+0.2\pi_{\st R}=\pi_{\st S},\qquad
      0.2\pi_{\st S}+0.4\pi_{\st C}+0.4\pi_{\st R}=\pi_{\st C},\qquad
      \pi_{\st S}+\pi_{\st C}+\pi_{\st R}=1 .
    \]
    The first gives $3\pi_{\st S}=3\pi_{\st C}+2\pi_{\st R}$, the second $6\pi_{\st C}=2\pi_{\st S}+4\pi_{\st R}$; solving,
    \[
      \pi=\big(\tfrac6{13},\ \tfrac4{13},\ \tfrac3{13}\big)\approx(0.4615,\ 0.3077,\ 0.2308),
    \]
    which agrees with the common row of $\PP^{20}$ on slide~\ref{fr:q1}. About $46\%$ of days are sunny in the long run (the precise sense of ``long run'' is given on slide~\ref{fr:ergodic}).
  \end{exampleblock}
  % python3: V/C: pi_V*0.8 = pi_C*0.5, pi_V+pi_C=1 -> (5/13, 8/13)
  \begin{exampleblock}{Vowel/consonant}
    $0.2\pi_{\st V}+0.5\pi_{\st C}=\pi_{\st V}$, i.e.\ $0.8\pi_{\st V}=0.5\pi_{\st C}$, with $\pi_{\st V}+\pi_{\st C}=1$: $\pi=(\tfrac5{13},\tfrac8{13})$, the limit row of slide~\ref{fr:twobytwo} and the weights in the Entropy lecture's entropy rate.
  \end{exampleblock}
\end{frame}

% ---------------------------------------------------------------------------
\begin{frame}{Detailed Balance and Reversible Chains}
  \label{fr:db}
  \small
  \begin{block}{Theorem}
    If $\pi\colon S\to[0,1]$ sums to $1$ and satisfies the \textbf{detailed balance} equations
    \[
      \pi_i\,p(i,j)=\pi_j\,p(j,i)\qquad\text{for all } i,j\in S,
    \]
    then $\pi$ is stationary. Such a chain is called \textbf{reversible} (with respect to $\pi$).
  \end{block}
  \words{If the flow along each edge balances in both directions, then the total flow into each state equals the flow out of it.}
  \begin{block}{Proof}
    $\sum_i\pi_i\,p(i,j)=\sum_i\pi_j\,p(j,i)=\pi_j\sum_ip(j,i)=\pi_j$. \hfill$\square$
  \end{block}
  \begin{itemize}
    \item Stronger than stationarity, with $|S|^2$ equations, but each involves only two unknowns, so the system is solved edge by edge; if it has no solution, $\pi$ must be found from $\pi\PP=\pi$.
    \item Random walk on a graph: $\pi_i=d(i)/\sum_kd(k)$ satisfies $\pi_ip(i,j)=\frac{1}{\sum_kd(k)}=\pi_jp(j,i)$ on every edge. Cube and cycle: uniform.
    \item Started from $\pi$, a reversible chain has the same joint distributions when run backwards, since $\Prob(X_n=i,X_{n+1}=j)=\pi_ip(i,j)=\pi_jp(j,i)=\Prob(X_n=j,X_{n+1}=i)$.
  \end{itemize}
\end{frame}

% ---------------------------------------------------------------------------
\begin{frame}{Detailed Balance: The Ehrenfest Urn}
  \label{fr:ehrenfest}
  \small
  % python3 sympy: Ehrenfest N=4, pi=(1,4,6,4,1)/16: pi P == pi checked; pi_i p(i,i+1) = pi_{i+1} p(i+1,i):
  %   1/16*1=4/16*1/4, 4/16*3/4=6/16*1/2, 6/16*1/2=4/16*3/4, 4/16*1/4=1/16*1.  All equal.
  \begin{exampleblock}{The Ehrenfest urn: $\pi$ is binomial}
    Transitions occur only between neighbouring states, so detailed balance is one equation for each edge: $\pi_i\frac{N-i}{N}=\pi_{i+1}\frac{i+1}{N}$, i.e.\ $\pi_{i+1}=\pi_i\frac{N-i}{i+1}$. Iterating from $\pi_0$ gives $\pi_i=\pi_0\binom Ni$, and $\sum_i\pi_i=1$ forces $\pi_0=2^{-N}$:
    \[
      \pi_i=\binom Ni2^{-N}\qquad\text{(each ball in urn 1 independently with probability } \tfrac12\text{)}.
    \]
    $N=4$: $\pi=\tfrac1{16}(1,4,6,4,1)$. Check one edge: $\pi_1p(1,2)=\tfrac4{16}\cdot\tfrac34=\tfrac3{16}=\tfrac6{16}\cdot\tfrac12=\pi_2p(2,1)$.
  \end{exampleblock}
  \words{Under the stationary distribution each ball is in urn~$1$ independently with probability $\tfrac12$.}
  \begin{itemize}
    \item Solving $\pi\PP=\pi$ directly means solving a $5\times5$ system, or a $101\times101$ system for $N=100$; detailed balance gives $\pi$ by a recursion along the edges.
    \item $\pi_0=2^{-N}$, so by slide~\ref{fr:recclass} the mean time between two moments when urn~$1$ is empty is $2^N$ steps: $16$ for $N=4$, about $1.3\cdot10^{30}$ for $N=100$; meanwhile $\pi$ is concentrated within a few multiples of $\sqrt N$ of $N/2$.
    % python3: 2**100 = 1.2677e30
  \end{itemize}
\end{frame}

% ---------------------------------------------------------------------------
\begin{frame}{PageRank: The Stationary Distribution of a Random Surfer}
  \label{fr:pagerank}
  \small
  \begin{columns}[T]
    \begin{column}{0.6\textwidth}
      \begin{block}{Definition}
        $n$ web pages, links as directed edges. A \textbf{random surfer} moves as follows: with probability $d$ (the damping factor, $d=0.85$) follow a uniformly chosen link out of the current page; with probability $1-d$ jump to a uniformly random page. Its transition matrix is
        \[
          \PP=d\,\mathbf H+\frac{1-d}{n}\,\mathbf J,\qquad \mathbf H_{ij}=\frac{\ind\{i\to j\}}{\#\{\text{links out of } i\}},\quad \mathbf J_{ij}=1 .
        \]
        The \textbf{PageRank} of page $j$ is $\pi_j$, where $\pi\colon S\to[0,1]$ is the stationary distribution of $\PP$.
      \end{block}
      \words{A page ranks highly if highly ranked pages link to it; the equation $\pi_j=\sum_i\pi_i\,p(i,j)$ makes this precise.}
    \end{column}
    \begin{column}{0.37\textwidth}
      \centering
      \begin{tikzpicture}[>=Stealth,font=\small,every node/.style={circle,draw,minimum size=7mm,inner sep=0pt}]
        \node (A) at (0,1.6) {$\st A$}; \node (B) at (2,1.6) {$\st B$}; \node (C) at (2,0) {$\st C$}; \node (D) at (0,0) {$\st D$};
        \draw[->,accent,thick] (A) -- (B); \draw[->,accent,thick] (A) -- (C); \draw[->,accent,thick] (B) -- (C);
        \draw[->,accent,thick] (C) to[bend left=25] (A); \draw[->,accent,thick] (D) -- (C);
      \end{tikzpicture}
      \par\medskip
      {\footnotesize Four pages. Links $\st A\to\st B,\st C$; $\st B\to\st C$; $\st C\to\st A$; $\st D\to\st C$. $\PP$ adds a $0.0375$ chance of jumping to each page from each page.}
      \par\medskip
      {\footnotesize In $\mathbf H_{ij}$, $i$ is the page the surfer is on and $j$ the page linked to; row sums are $1$ (for a page with no out-links put $\mathbf H_{ij}=1/n$ for all $j$).}
    \end{column}
  \end{columns}
  \begin{itemize}
    \item The jump is needed because from a page with no out-links, or a group of pages with no exits, a surfer following links alone can never leave, and the chain with matrix $\mathbf H$ alone may have no unique $\pi$ or may fail to converge (slides~\ref{fr:classes}--\ref{fr:convwhy}). With $d<1$ every entry of $\PP$ is positive, so the results below apply.
  \end{itemize}
\end{frame}

% ---------------------------------------------------------------------------
\begin{frame}{PageRank: The Four-Page Web}
  \label{fr:pagerank2}
  \small
  % python3 numpy/sympy: links A->B,C; B->C; C->A; D->C; d=0.85:
  %   G = 0.85 H + 0.0375 J = [[.0375,.4625,.4625,.0375],[.0375,.0375,.8875,.0375],[.8875,.0375,.0375,.0375],[.0375,.0375,.8875,.0375]]
  %   pi = (659/1769, 27713/141520, 2789/7076, 3/80) = (0.3725, 0.1958, 0.3941, 0.0375); equals G^30 row 0.
  %   without damping H^30 row 0 = (0.4, 0.2, 0.4, 0).
  \[
    \mathbf H=\begin{pmatrix}0&\tfrac12&\tfrac12&0\\0&0&1&0\\1&0&0&0\\0&0&1&0\end{pmatrix},\qquad
    \PP=0.85\,\mathbf H+0.0375\,\mathbf J=\begin{pmatrix}0.0375&0.4625&0.4625&0.0375\\0.0375&0.0375&0.8875&0.0375\\0.8875&0.0375&0.0375&0.0375\\0.0375&0.0375&0.8875&0.0375\end{pmatrix}.
  \]
  \begin{exampleblock}{Solving $\pi\PP=\pi$, $\sum_j\pi_j=1$}
    \[
      \pi\approx(0.373,\ 0.196,\ 0.394,\ 0.038)\quad\text{for }(\st A,\st B,\st C,\st D),\qquad\text{exactly }\pi_{\st D}=\tfrac{3}{80}=0.0375 .
    \]
    $\st C$ ranks first; three pages link to it, and one of them, $\st A$, is itself highly ranked. $\st A$ is second although it has a single in-link, because that link comes from the top page and is $\st C$'s only link. $\st D$ has no in-links and receives only the jump mass $\frac{1-d}{n}$.
  \end{exampleblock}
  \begin{itemize}
    \item The definition of rank is circular, since $\st A$ ranks highly because $\st C$ does and $\st C$ because $\st A$ does. The stationary equation expresses this circularity as a fixed-point condition, and the powers $\PP^n$ compute the fixed point by iteration ($\PP^{30}$ has all rows equal to $\pi$ to four places).
    \item Without damping ($d=1$), $\pi=(0.4,0.2,0.4,0)$ is stationary for $\mathbf H$: a page with no in-links has rank $0$.
  \end{itemize}
\end{frame}


% ===========================================================================
\begin{frame}{Communicating Classes and Irreducibility}
  \label{fr:classes}
  \small
  \begin{block}{Definition}
    State $i$ \textbf{leads to} $j$, written $i\to j$, if $p^{(n)}(i,j)>0$ for some $n\ge0$. States $i,j$ \textbf{communicate}, $i\leftrightarrow j$, if $i\to j$ and $j\to i$. This is an equivalence relation on $S$; its equivalence classes are the \textbf{communicating classes}. A class $C$ is \textbf{closed} if $i\in C$ and $i\to j$ imply $j\in C$. The chain is \textbf{irreducible} if $S$ is a single class.
  \end{block}
  \words{Two states are in the same class if the chain can get from each to the other; a closed class is one the chain cannot leave; irreducible means every state is reachable from every other.}
  \begin{itemize}
    \item $i\to j$ iff there is a path $i=i_0,i_1,\dots,i_n=j$ with every $p(i_{k-1},i_k)>0$: a directed path in the graph whose edges are the positive entries of $\PP$. The classes are the strongly connected components of that graph.
    \item Weather, vowel/consonant, Ehrenfest, cycle, cube, PageRank with $d<1$: irreducible.
    \item Gambler's ruin: three classes, $\{0\}$ and $\{N\}$ closed (absorbing states), $\{1,\dots,N-1\}$ not closed. Four-page web without damping: $\{\st A,\st B,\st C\}$ closed, $\{\st D\}$ not closed.
    \item A state in a class that is not closed is left permanently with positive probability at each visit, so it is visited only finitely often (next slides).
  \end{itemize}
\end{frame}

% ---------------------------------------------------------------------------
\begin{frame}{Periodicity}
  \label{fr:period}
  \small
  \begin{block}{Definition}
    The \textbf{period} of state $i$ is $d(i)=\gcd\{n\ge1 : p^{(n)}(i,i)>0\}$ (with $d(i)=\infty$ if the set is empty). State $i$ is \textbf{aperiodic} if $d(i)=1$. A chain is aperiodic if every state is.
  \end{block}
  \words{The period is the greatest common divisor of the lengths of all round trips from $i$; period $2$ means every return takes an even number of steps.}
  \begin{block}{Theorem}
    All states in a communicating class have the same period. If $p(i,i)>0$ for some $i$ in a class, the class is aperiodic.
  \end{block}
  \begin{block}{Proof}
    Let $i\leftrightarrow j$ with $p^{(a)}(i,j)>0$, $p^{(b)}(j,i)>0$. Every round trip from $j$ of length $n$ gives round trips from $i$ of lengths $a+b$ and $a+n+b$, so $d(i)$ divides both, hence divides $n$; thus $d(i)\le d(j)$, and symmetrically. If $p(i,i)>0$ then $1$ is a round-trip length. \hfill$\square$
  \end{block}
\end{frame}

% ---------------------------------------------------------------------------
\begin{frame}{Periodicity: The Random Walk on the Cube}
  \label{fr:cubeparity}
  \small
  % python3 numpy: cube P^10 row 000 = (.25,0,0,.25,0,.25,.25,0); P^11 row 000 = (0,.25,.25,0,.25,0,0,.25); P^n(000,000)=0 for odd n
  \begin{columns}[T]
    \begin{column}{0.63\textwidth}
      \begin{exampleblock}{Random walk on the cube}
        Each step flips one coordinate, so the parity of the number of $1$s alternates: even at even times, odd at odd times. Hence $p^{(n)}(000,000)=0$ for odd $n$, $d=2$, and
        \begin{align*}
          (\PP^{10})_{000,\cdot}&\approx\tfrac14(1,0,0,1,0,1,1,0),\\
          (\PP^{11})_{000,\cdot}&\approx\tfrac14(0,1,1,0,1,0,0,1)
        \end{align*}
        (order $000,100,010,110,001,101,011,111$; zeros exact; $p^{(10)}(000,000)=\tfrac14+\tfrac34\,3^{-10}$). The rows of $\PP^n$ do not converge: they alternate between two limits, while $\pi=(\tfrac18,\dots,\tfrac18)$ is stationary all along.
      \end{exampleblock}
      \words{Started at $000$, the parity of the number of $1$s in $X_n$ equals the parity of $n$; $\pi$ exists, but $X_n$ does not converge in distribution.}
    \end{column}
    \begin{column}{0.34\textwidth}
      \centering
      \begin{tikzpicture}[scale=0.85,font=\scriptsize]
        \coordinate (000) at (0,0); \coordinate (100) at (1.4,0); \coordinate (010) at (0,1.4); \coordinate (110) at (1.4,1.4);
        \coordinate (001) at (0.6,0.5); \coordinate (101) at (2.0,0.5); \coordinate (011) at (0.6,1.9); \coordinate (111) at (2.0,1.9);
        \draw[soft] (000)--(100)--(110)--(010)--cycle; \draw[soft] (001)--(101)--(111)--(011)--cycle;
        \draw[soft] (000)--(001) (100)--(101) (110)--(111) (010)--(011);
        \foreach \v in {000,110,101,011}{\fill[accent] (\v) circle (2.4pt);}
        \foreach \v in {100,010,001,111}{\fill[amber] (\v) circle (2.4pt);}
        \node[below left] at (000) {$000$}; \node[above right] at (111) {$111$};
      \end{tikzpicture}
      \par\smallskip
      {\footnotesize Blue: even parity; amber: odd. Every edge joins the two colours, so the walk alternates colours forever.}
    \end{column}
  \end{columns}
  \begin{itemize}
    \item \textbf{Cycle} of length $L$: $d=2$ if $L$ is even, $d=1$ if $L$ is odd (round trips of lengths $2$ and $L$). \textbf{Ehrenfest}: $d=2$. \textbf{Weather}, \textbf{vowel/consonant}: $p(i,i)>0$, so $d=1$.
    \item The lazy chain $\PP'=\tfrac12(\mathbf I+\PP)$ has the same $\pi$ and period $1$, so it removes periodicity.
  \end{itemize}
\end{frame}

% ---------------------------------------------------------------------------

% ---------------------------------------------------------------------------
\begin{frame}{Recurrence and Transience}
  \label{fr:recurrence}
  \small
  \begin{block}{Definition}
    The \textbf{return time} to $i$ is $T_i\colon\Omega\to\{1,2,\dots\}\cup\{\infty\}$, $T_i=\min\{n\ge1 : X_n=i\}$ ($\infty$ if there is no such $n$). State $i$ is \textbf{recurrent} if $\Prob_i(T_i<\infty)=1$ and \textbf{transient} if $\Prob_i(T_i<\infty)<1$. A recurrent state is \textbf{positive recurrent} if $\E_i[T_i]<\infty$ and \textbf{null recurrent} otherwise.
  \end{block}
  \words{A state is recurrent if the chain started there returns with probability one, and transient if with positive probability it never returns.}
  \begin{itemize}
    \item $T_i$ is a row quantity, the first time after $0$ at which the path is in state $i$. It is a waiting time of the kind studied in the Conditional Expectation lecture; for the gambler, the duration $T$ of the game was such a time.
    \item Gambler's ruin: $0$ and $N$ are recurrent ($T_i=1$ surely); $1,\dots,N-1$ are transient, since the game ends with probability one and the chain then never returns.
    \item Random walk on $\Z$: state $0$ is recurrent iff $p=\tfrac12$ (slide~\ref{fr:walkZ}). For $p=\tfrac12$, $0$ is recurrent but $\E_0[T_0]=\infty$, so the fair walk is null recurrent; this is the statement ``ruin certain but $\E[T]=\infty$'' of slide~\ref{fr:regimes}.
    \item For a positive recurrent state, returns occur every $\E_i[T_i]$ steps on average, so the long-run fraction of time spent at $i$ is $1/\E_i[T_i]$ (slide~\ref{fr:ergodic}).
  \end{itemize}
\end{frame}

% ---------------------------------------------------------------------------
\begin{frame}{Recurrence: The Visits Criterion}
  \label{fr:visits}
  \small
  \begin{block}{Theorem}
    Let $V_i\colon\Omega\to\{0,1,\dots\}\cup\{\infty\}$, $V_i=\sum_{n\ge1}\ind\{X_n=i\}$, be the number of returns to $i$, and $f=\Prob_i(T_i<\infty)$. Then
    \[
      \Prob_i(V_i\ge k)=f^{\,k},\qquad \E_i[V_i]=\sum_{n\ge1}p^{(n)}(i,i)=\frac{f}{1-f}\quad(=\infty\text{ if } f=1).
    \]
    Hence $i$ is recurrent iff $\sum_n p^{(n)}(i,i)=\infty$ iff $\Prob_i(V_i=\infty)=1$; transient iff $\sum_n p^{(n)}(i,i)<\infty$ iff $V_i<\infty$ with probability one.
  \end{block}
  \words{The number of returns to $i$ is geometric with parameter $f$, so it is infinite with probability one in the recurrent case and has finite mean in the transient case.}
  \begin{block}{Proof}
    Each time the chain is at $i$ it restarts (Markov property), and returns once more with probability $f$, independently of the past; so $V_i\ge k$ requires $k$ successive returns: $f^k$. Then $\E_i[V_i]=\sum_{k\ge1}\Prob_i(V_i\ge k)=\sum_kf^k$, and also $\E_i[V_i]=\sum_n\E_i[\ind\{X_n=i\}]=\sum_np^{(n)}(i,i)$ by linearity, column by column. \hfill$\square$
  \end{block}
\end{frame}

% ---------------------------------------------------------------------------
\begin{frame}{Recurrence: Class Property and Finite Chains}
  \label{fr:recclass}
  \small
  \begin{block}{Theorem}
    \begin{enumerate}
      \item[(a)] Recurrence and transience are class properties: all states of a class are recurrent, or all are transient.
      \item[(b)] Every finite closed class is recurrent; every class that is not closed is transient. In particular a finite irreducible chain is recurrent, indeed positive recurrent, and its stationary distribution is unique with $\pi_i=1/\E_i[T_i]$.
    \end{enumerate}
  \end{block}
  \words{In a finite chain a state is transient only if its class can be left permanently; within a finite closed class every state is revisited, on average every $1/\pi_i$ steps.}
  \begin{block}{Proof sketch}
    (a) If $i\leftrightarrow j$ with $p^{(a)}(i,j)p^{(b)}(j,i)>0$, then $p^{(a+n+b)}(i,i)\ge p^{(a)}(i,j)p^{(n)}(j,j)p^{(b)}(j,i)$, so the two series $\sum_np^{(n)}(i,i)$, $\sum_np^{(n)}(j,j)$ diverge together. (b) Started in a finite closed class, some state is visited infinitely often, so that state is recurrent, so all are by (a). A class that is not closed can be left permanently. The formula $\pi_i=1/\E_i[T_i]$ is the ergodic theorem, slide~\ref{fr:ergodic}. \hfill$\square$
  \end{block}
  % python3: 1/pi_R = 13/3 = 4.333; 1/pi_0 (Ehrenfest N=4) = 16; V/C mean return to V = 1+0.8*2 = 13/5 = 1/pi_V
  \begin{itemize}
    \item Weather: mean return time to $\st R$ is $1/\pi_{\st R}=\tfrac{13}3$ days. Vowel/consonant: mean return to $\st V$ is $1/\pi_{\st V}=\tfrac{13}5$, in agreement with first-step analysis: from $\st C$ the wait for $\st V$ is geometric with mean $2$, giving $1+0.8\cdot2$.
  \end{itemize}
\end{frame}

% ---------------------------------------------------------------------------
\begin{frame}{The Random Walk on $\Z$ Is Recurrent If and Only If It Is Fair}
  \label{fr:walkZ}
  \small
  \begin{block}{Theorem}
    The simple random walk on $\Z$ is recurrent if $p=\tfrac12$ (null recurrent) and transient if $p\ne\tfrac12$.
  \end{block}
  \words{Without drift the walk returns to $0$ infinitely often, with infinite mean return time; with drift it returns only finitely often.}
  % python3: C(2n,n)/4^n vs 1/sqrt(pi n): n=1: 0.5 vs 0.564; n=5: 0.246 vs 0.252; n=50: 0.0796 vs 0.0798; 4pq at p=0.6 is 0.96
  \begin{block}{Proof}
    Returns to $0$ take an even number of steps, $n$ up and $n$ down: $p^{(2n)}(0,0)=\binom{2n}{n}p^nq^n$ (slide~\ref{fr:distSn}). By Stirling's formula (not proved here) $\binom{2n}n\approx4^n/\sqrt{\pi n}$, the ratio tending to $1$, so
    \[
      p^{(2n)}(0,0)\approx\frac{(4pq)^n}{\sqrt{\pi n}} .
    \]
    For $p=\tfrac12$, $4pq=1$ and $\sum_n(\pi n)^{-1/2}=\infty$, so $0$ is recurrent by the visits criterion. For $p\ne\tfrac12$, $4pq=1-(p-q)^2<1$ and the series converges, so $0$ is transient. Null recurrence: $\E_0[T_0]=\infty$ by slide~\ref{fr:regimes}. \hfill$\square$
  \end{block}
  \begin{itemize}
    % python3: sum_{n>=1} C(2n,n)(pq)^n at p=0.6 -> 4.000 (partial sum to n=4000); 1/sqrt(1-0.96)-1 = 4; f = 4/5 = 1-|p-q|
    \item $p=0.6$: $4pq=0.96$ and $\E_0[V_0]=\sum_{n\ge1}\binom{2n}{n}(pq)^n=(1-4pq)^{-1/2}-1=4$, so $f=\tfrac45=1-|p-q|$ by the visits criterion; the walk returns four times on average.
  \end{itemize}
\end{frame}


% ===========================================================================
\begin{frame}{Hitting Probabilities: The Minimal Solution of a First-Step System}
  \label{fr:hitting}
  \small
  \begin{block}{Definition}
    For $A\subseteq S$, the \textbf{hitting time} of $A$ is $T_A\colon\Omega\to\{0,1,2,\dots\}\cup\{\infty\}$, $T_A=\min\{n\ge0 : X_n\in A\}$, and the \textbf{hitting probability} is the function
    \[
      h^A\colon S\to[0,1],\qquad h^A_i=\Prob_i(T_A<\infty).
    \]
  \end{block}
  \begin{block}{Theorem}
    $h^A$ is the \textbf{minimal non-negative solution} of the system
    \[
      h_i=1\ \ (i\in A),\qquad h_i=\sum_{j\in S}p(i,j)\,h_j\ \ (i\notin A):
    \]
    it solves it, and every solution $x$ with all $x_i\ge0$ has $x_i\ge h^A_i$ for all $i$.
  \end{block}
  \words{Off the target set, the hitting probability from $i$ is the $p(i,\cdot)$-weighted average of the hitting probabilities from the next states; if the system has several non-negative solutions, the hitting probability is the smallest.}
  \begin{itemize}
    \item $A$ is the fixed target set; $i$ is the starting state, the variable of $h^A$. Gambler's ruin: $S=\{0,\dots,N\}$, $A=\{N\}$ (next slide); the exit problem of slide~\ref{fr:firststep}: $A=\{-a,a\}$.
  \end{itemize}
\end{frame}

% ---------------------------------------------------------------------------
\begin{frame}{Hitting Probabilities: Proof}
  \label{fr:hittingproof}
  \small
  \begin{block}{Proof}
    Solves it: $h^A_i=1$ on $A$; off $A$, condition on the first step (total probability over the rows $\{X_1=j\}$) and use the Markov property: $\Prob_i(T_A<\infty)=\sum_jp(i,j)\Prob_j(T_A<\infty)$.
    Minimal: let $x\ge0$ solve the system. Then $x_i\ge\Prob_i(T_A\le n)$ for all $n$, by induction: true for $n=0$ (on $A$, $x_i=1$; off $A$, $x_i\ge0=\Prob_i(T_A\le0)$); and for $i\notin A$,
    \[
      x_i=\sum_jp(i,j)\,x_j\ \ge\ \sum_jp(i,j)\,\Prob_j(T_A\le n)=\Prob_i(T_A\le n+1)
    \]
    by the induction hypothesis and first-step analysis. Let $n\to\infty$: $x_i\ge\Prob_i(T_A<\infty)=h^A_i$. \hfill$\square$
  \end{block}
  \words{Any non-negative solution dominates the probability of hitting within $n$ steps for every $n$, hence dominates the probability of ever hitting.}
  \begin{itemize}
    \item On a finite state space in which $A$ can be reached from every state, the system has a unique solution and minimality plays no role; when there are absorbing states outside $A$, or the state space is infinite, minimality selects the correct solution.
  \end{itemize}
\end{frame}

% ---------------------------------------------------------------------------
\begin{frame}{Hitting Probabilities: The Gambler's Ruin Is the Case $A=\{N\}$}
  \label{fr:hittingruin}
  \small
  \begin{exampleblock}{Finite ruin chain, $A=\{N\}$}
    The system reads $h_N=1$; $h_i=p\,h_{i+1}+q\,h_{i-1}$ for $0<i<N$; and at the absorbing state $0$: $h_0=\sum_jp(0,j)h_j=h_0$, which imposes no constraint. The general solution is $h_i=a+b\,r^i$ (or $a+bi$ if $p=q$) for $i\ge1$ with $h_0$ free; non-negative solutions exist with any $h_0\in[0,\ 1]$, and the minimal one has $h_0=0$, which recovers slides~\ref{fr:ruinfair}--\ref{fr:ruinunfair}.
  \end{exampleblock}
  \words{The boundary condition $h_0=0$ imposed on slide~\ref{fr:ruineq} is supplied here by minimality; an absorbing state has hitting probability $0$ for any target set not containing it.}
  \begin{itemize}
    \item The hitting probabilities of the ruin problem come from this system with different targets: $A=\{N\}$ gives the probability of reaching the target, $A=\{0\}$ the probability of ruin, and $A=\{0,N\}$ gives $h\equiv1$, so the game ends with probability one.
    \item The ruin argument of slide~\ref{fr:ruineq}, conditioning on the first round and restarting, is the proof that $h^A$ solves the system; the only addition for general chains is minimality.
  \end{itemize}
\end{frame}

% ---------------------------------------------------------------------------
\begin{frame}{Why ``Minimal'': The Infinite Ruin Chain}
  \label{fr:minimal}
  \small
  \begin{exampleblock}{Infinite ruin chain, $S=\{0,1,2,\dots\}$, $A=\{0\}$ (slide~\ref{fr:infinite} again)}
    $h_0=1$ and $h_i=p\,h_{i+1}+q\,h_{i-1}$ for $i\ge1$; general solution $h_i=1+b\,(r^i-1)$ with $r=q/p$.
    \begin{itemize}
      \item $p>q$ ($r<1$): $h_i\ge0$ for all $i$ requires $b\le1$, and $h_i$ decreases as $b$ increases, so the minimal solution is $b=1$: $h_i=r^i=(q/p)^i$.
      \item $p<q$ ($r>1$): $h_i\ge0$ for all $i$ requires $b\ge0$ (otherwise $h_i\to-\infty$), minimal at $b=0$: $h_i=1$. \quad $p=q$: $h_i=1+bi$, $b\ge0$, minimal $h_i=1$.
    \end{itemize}
    $h\equiv1$ solves the system in every case; it is the answer only when it is the smallest non-negative solution.
  \end{exampleblock}
  \words{The system alone does not distinguish certain ruin from ruin with probability $(q/p)^i$; the minimality condition does, and the result agrees with the limit $N\to\infty$ on slide~\ref{fr:infinite}.}
  \begin{itemize}
    \item The same occurs for transient chains in general: on an infinite state space, a non-negative solution can be inflated by adding a positive solution of the homogeneous equation that vanishes on $A$; only the smallest is the hitting probability.
  \end{itemize}
\end{frame}

% ---------------------------------------------------------------------------
\begin{frame}{Expected Hitting Times}
  \label{fr:khit}
  \small
  \begin{block}{Theorem}
    Let $k^A\colon S\to[0,\infty]$, $k^A_i=\E_i[T_A]$. Then $k^A$ is the minimal non-negative solution of
    \[
      k_i=0\ \ (i\in A),\qquad k_i=1+\sum_{j\in S}p(i,j)\,k_j\ \ (i\notin A).
    \]
  \end{block}
  \words{Off the target set, the expected hitting time is one step plus the weighted average of the expected hitting times from the next states; the exit-time and duration equations of slides~\ref{fr:firststep} and~\ref{fr:duration} are instances of this system.}
  \begin{block}{Proof}
    Solves it: off $A$, $T_A=1+T_A'$ where $T_A'$ is the hitting time of the chain restarted at $X_1$; take $\E_i$ and condition on $X_1$. Minimal: as before, any non-negative solution $x$ satisfies $x_i\ge\E_i[\min(T_A,n)]$ for all $n$, by induction. \hfill$\square$
  \end{block}
  \begin{itemize}
    \item The system may have no finite non-negative solution; then $k^A_i=\infty$ for some $i$, as for the fair walk on $\Z$. A finite solution, when the target is reachable from everywhere in a finite chain, is unique.
    \item In matrix form on a finite chain: $(\mathbf I-\mathbf Q)\,k=\mathbf 1$ with $\mathbf Q$ the transitions among the states off $A$; slide~\ref{fr:absorbing} inverts it.
  \end{itemize}
\end{frame}

% ---------------------------------------------------------------------------
\begin{frame}{Expected Hitting Times: Corner to Opposite Corner on the Cube}
  \label{fr:cubehit}
  \small
  We compute the expected number of steps for the random walk on the cube to go from $000$ to $111$.
  \begin{block}{Direct computation}
    $\E_{000}[T_{111}]=\sum_nn\,\Prob_{000}(T_{111}=n)$, where each $\Prob_{000}(T_{111}=n)$ is a sum over paths of length $n$ that avoid $111$ until their last step; the quantity is not determined by the distribution of $X_n$ at any fixed time.
  \end{block}
  % python3 sympy: k3=0, k0=1+k1, k1=1+k0/3+2k2/3, k2=1+2k1/3 -> k0=10, k1=9, k2=7.  Adjacent target: 7, 9, 10.
  \begin{exampleblock}{Random walk on the cube: from $000$ to the opposite corner $111$}
    Group states by distance $d\in\{0,1,2,3\}$ from $111$; by symmetry $k$ depends only on $d$. From distance $d$ a step goes to distance $d-1$ with probability $d/3$ and to $d+1$ with probability $(3-d)/3$:
    \[
      k_0=0,\qquad k_1=1+\tfrac13k_0+\tfrac23k_2,\qquad k_2=1+\tfrac23k_1+\tfrac13k_3,\qquad k_3=1+k_2 ,
    \]
    giving $k_1=7$, $k_2=9$, $k_3=10$: ten steps on average from a corner to the opposite corner.
  \end{exampleblock}
  \words{Symmetry reduces the eight unknowns to three, and first-step analysis replaces the infinite sum over paths by three linear equations.}
  \begin{itemize}
    \item Check: from a neighbour of $000$ the expected time back to $000$ is $7$, so the mean return time to $000$ is $1+7=8=1/\pi_{000}$, as the ergodic theorem predicts (slide~\ref{fr:ergodic}).
  \end{itemize}
\end{frame}

% ---------------------------------------------------------------------------
\begin{frame}{Return to Question 3: Days Until Rain}
  \label{fr:q3}
  \small
  % python3 sympy: kS = 1 + 0.7 kS + 0.2 kC, kC = 1 + 0.3 kS + 0.4 kC -> kS = 20/3, kC = 5.
  %   mean return time to R: 1 + 0.2*20/3 + 0.4*5 = 13/3 = 1/pi_R.  Checked.
  $A=\{\st R\}$, $k_{\st R}=0$, and for the other two states
  \[
    k_{\st S}=1+0.7\,k_{\st S}+0.2\,k_{\st C},\qquad k_{\st C}=1+0.3\,k_{\st S}+0.4\,k_{\st C}.
  \]
  \begin{exampleblock}{Answer}
    From the second, $0.6\,k_{\st C}=1+0.3\,k_{\st S}$; substituting into the first, $0.3\,k_{\st S}=1+0.2\,k_{\st C}$, so $k_{\st S}=\tfrac{20}3\approx6.7$ days and $k_{\st C}=5$ days.
  \end{exampleblock}
  \begin{block}{Consistency check with the stationary distribution}
    The mean return time to $\st R$ is one step plus the expected hitting time from wherever that step lands:
    \[
      \E_{\st R}[T_{\st R}]=1+0.2\,k_{\st S}+0.4\,k_{\st C}+0.4\cdot0=1+\tfrac43+2=\tfrac{13}3=\frac1{\pi_{\st R}} .
    \]
  \end{block}
  \words{The three opening questions have been answered by a matrix power, a left eigenvector, and a first-step linear system respectively.}
  \begin{itemize}
    \item For the vowel/consonant chain, the expected number of letters until the first vowel is $2$ from $\st C$ (a geometric waiting time) and $1+0.8\cdot2=2.6$ from $\st V$; until the first occurrence of $\st{VV}$ it is $k_{\st C}=15$ and $k_{\st V}=13$, from the first-step system $k_{\st C}=2+k_{\st V}$, $k_{\st V}=1+0.8\,k_{\st C}$.
    % python3 sympy: kC = 2 + kV, kV = 1 + 0.8 kC -> kV = 13, kC = 15; checked via fundamental matrix of the 3-state pattern chain.
  \end{itemize}
\end{frame}


% ===========================================================================
\begin{frame}{The Convergence Theorem}
  \label{fr:conv}
  \small
  \begin{block}{Theorem}
    Let $\PP$ be irreducible and aperiodic on a finite state space $S$. Then $\PP$ has a unique stationary distribution $\pi$, and for all $i,j\in S$,
    \[
      p^{(n)}(i,j)\ \xrightarrow{\ n\to\infty\ }\ \pi_j .
    \]
    Equivalently, for every initial distribution $\lambda$, $\lambda\PP^n\to\pi$: every row of $\PP^n$ converges to $\pi$.
  \end{block}
  \words{For large $n$ the distribution of $X_n$ is close to $\pi$, whatever the initial distribution.}
  \begin{itemize}
    \item Slide~\ref{fr:twobytwo}: rows of the vowel/consonant $\PP^n$ approaching $(\tfrac5{13},\tfrac8{13})$. Slide~\ref{fr:q1}: rows of the weather $\PP^7$ agreeing to two places. PageRank: $\PP^{30}$ has all rows equal to $\pi$ to four places, which is how $\pi$ is computed for $10^{10}$ pages.
    \item The theorem concerns the \emph{columns} of the table of paths, the distribution at a fixed late time. It does not describe an individual path, which continues to move; the ergodic theorem (slide~\ref{fr:ergodic}) is the corresponding statement about rows.
    \item Both hypotheses are necessary: irreducibility for the uniqueness of $\pi$ and the independence of the limit from the start, aperiodicity for the convergence itself. Examples follow on the next slide.
  \end{itemize}
\end{frame}

% ---------------------------------------------------------------------------
\begin{frame}{Why Each Hypothesis Is Needed}
  \label{fr:convwhy}
  \small
  \begin{block}{Reducible: the limit depends on the start}
    Two closed classes, $\PP=\begin{pmatrix}\PP_1&0\\0&\PP_2\end{pmatrix}$. Rows of $\PP^n$ starting in class $1$ converge to $(\pi_1,0)$, in class $2$ to $(0,\pi_2)$; every mixture $(\alpha\pi_1,(1-\alpha)\pi_2)$ is stationary. Gambler's ruin: rows converge to $(1-h_i,0,\dots,0,h_i)$, different for each $i$; for $N=4$, $p=\tfrac12$, row $i$ tends to $(1-\tfrac i4,0,0,0,\tfrac i4)$.
  \end{block}
  \begin{block}{Periodic: the rows do not converge}
    The cube (slide~\ref{fr:cubeparity}): rows of $\PP^n$ alternate between two vectors, though $\pi=(\tfrac18,\dots,\tfrac18)$ is unique. Likewise the Ehrenfest urn and the cycle of even length.
  \end{block}
  \begin{block}{Infinite $S$: no stationary distribution}
    The fair walk on $\Z$ is irreducible, but has no stationary distribution (slide~\ref{fr:statexist}) and $p^{(2n)}(0,0)\approx1/\sqrt{\pi n}\to0$ (slide~\ref{fr:walkZ}): every entry of every row of $\PP^n$ tends to $0$, and the mass spreads out over ever more states.
  \end{block}
  \words{Without irreducibility the limit depends on the starting state, without aperiodicity there may be no limit, and on an infinite state space the limit may be zero at every state.}
\end{frame}

% ---------------------------------------------------------------------------
\begin{frame}{How Fast: The Second Eigenvalue}
  \label{fr:convrate}
  \small
  \begin{block}{Theorem (rate of convergence, diagonalisable case)}
    Suppose $\PP$ is diagonalisable, with eigenvalues $1=\lambda_1,\lambda_2,\dots,\lambda_{|S|}$, right eigenvectors $\mathbf v_k$ (columns, $\PP\mathbf v_k=\lambda_k\mathbf v_k$) and left eigenvectors $\mathbf u_k$ (rows, $\mathbf u_k\PP=\lambda_k\mathbf u_k$), chosen so that $\mathbf u_k\mathbf v_k=1$ and $\mathbf u_k\mathbf v_l=0$ for $k\ne l$. Then $\mathbf v_1=\mathbf 1$, $\mathbf u_1=\pi$, $\PP=\sum_k\lambda_k\mathbf v_k\mathbf u_k$, and
    \[
      \PP^n=\mathbf 1\pi+\sum_{k\ge2}\lambda_k^n\,\mathbf v_k\mathbf u_k,\qquad
      \big|p^{(n)}(i,j)-\pi_j\big|\le C\,|\lambda_2|^n,\quad |\lambda_2|=\max_{k\ge2}|\lambda_k|<1 .
    \]
  \end{block}
  \words{Each term of $\PP^n$ other than $\mathbf 1\pi$ carries a factor $\lambda_k^n$ with $|\lambda_k|<1$, so the difference from the limit decreases geometrically at a rate set by the second-largest eigenvalue modulus.}
  \begin{block}{Proof}
    $\PP\mathbf 1=\mathbf 1$ (row sums) and $\pi\PP=\pi$, $\pi\mathbf 1=1$ give $\mathbf v_1=\mathbf 1$, $\mathbf u_1=\pi$. As $(\mathbf v_k\mathbf u_k)(\mathbf v_l\mathbf u_l)=(\mathbf u_k\mathbf v_l)\,\mathbf v_k\mathbf u_l$ vanishes for $k\ne l$, powers act termwise: $\big(\sum_k\lambda_k\mathbf v_k\mathbf u_k\big)^n=\sum_k\lambda_k^n\mathbf v_k\mathbf u_k$. Irreducibility rules out a second eigenvalue $1$ (a second stationary distribution), aperiodicity the other eigenvalues of modulus $1$ (the cube has $-1$); these two facts are not proved here. \hfill$\square$
  \end{block}
  % python3 numpy: weather eigenvalues 1, 0.4562, 0.0438; max|P^n - pi|: n=1 .262, 3 .056, 6 .0053, 7 .0024.
  %   V/C eigenvalues 1, -0.3; max|P^n-pi| = 8/13 0.3^n: n=5 .0015, n=6 .00045.
  \begin{itemize}
    \item Vowel/consonant: $\lambda_2=-0.3$, error $\tfrac8{13}(0.3)^n$ (slide~\ref{fr:twobytwo}). Weather: eigenvalues $1,\ 0.456,\ 0.044$; $\max_{i,j}|p^{(n)}(i,j)-\pi_j|=0.26,\ 0.056,\ 0.0053,\ 0.0024$ at $n=1,3,6,7$.
  \end{itemize}
\end{frame}

% ---------------------------------------------------------------------------
\begin{frame}{Proof Sketch by Coupling}
  \label{fr:coupling}
  \small
  \begin{block}{Idea}
    Run two independent copies of the chain: $X$ started at $i$, and $Y$ started from $\pi$, so that $Y_n\sim\pi$ for every $n$ (slide~\ref{fr:statdef}). Let $\tau=\min\{n : X_n=Y_n\}$ be the first time they meet. Define $Z_n=X_n$ for $n<\tau$ and $Z_n=Y_n$ for $n\ge\tau$: switch to following $Y$ once they coincide.
  \end{block}
  \begin{block}{Two facts}
    \begin{enumerate}
      \item $Z$ is again a Markov chain with matrix $\PP$ started at $i$ (at the switch, both copies are at the same state and both continue by the same rule), so $\Prob(Z_n=j)=p^{(n)}(i,j)$.
      \item $Z_n=Y_n$ whenever $\tau\le n$, so $\big|p^{(n)}(i,j)-\pi_j\big|=\big|\Prob(Z_n=j)-\Prob(Y_n=j)\big|\le\Prob(\tau>n)$.
    \end{enumerate}
  \end{block}
  \begin{block}{Where the hypotheses enter}
    The pair $(X_n,Y_n)$ is a Markov chain on $S\times S$ with transitions $p(i,i')p(j,j')$. Irreducibility and aperiodicity of $\PP$ make this pair chain irreducible (aperiodicity is needed: for the cube, a pair at different parities can never meet), hence recurrent on the finite $S\times S$, so it reaches the diagonal: $\Prob(\tau<\infty)=1$, i.e.\ $\Prob(\tau>n)\to0$. \hfill$\square$
  \end{block}
  \words{After the two chains first meet, the chain started at $i$ can be replaced by the chain started from $\pi$ without changing its distribution, and the two are certain to meet.}
\end{frame}

% ---------------------------------------------------------------------------
\begin{frame}{The Ergodic Theorem: Long-Run Fractions of Time}
  \label{fr:ergodic}
  \small
  \begin{block}{Theorem}
    Let $\PP$ be irreducible on a finite $S$ with stationary distribution $\pi$ (no aperiodicity required). Let $V_i(n)\colon\Omega\to\{0,\dots,n\}$, $V_i(n)=\sum_{k=0}^{n-1}\ind\{X_k=i\}$, be the number of visits to $i$ before time $n$. Unlike $V_i$ of slide~\ref{fr:visits}, this count includes time $0$, so $V_i(n)\to V_i+\ind\{X_0=i\}$. Then, whatever the initial distribution, with probability one
    \[
      \frac{V_i(n)}{n}\ \xrightarrow{\ n\to\infty\ }\ \pi_i\quad\text{for every } i,
      \qquad\text{and}\qquad
      \frac1n\sum_{k=0}^{n-1}f(X_k)\ \to\ \sum_{i\in S}\pi_i\,f(i)\quad\text{for every } f\colon S\to\R .
    \]
    Moreover $\pi_i=1/\E_i[T_i]$.
  \end{block}
  \words{Along a single path, the fraction of time spent in state $i$ converges to $\pi_i$, and the time average of any function of the state converges to its average under $\pi$.}
  \begin{itemize}
    \item $V_i(n)/n$ is a \emph{row} quantity, computed by counting along one path. The theorem says that almost every row has the same limiting frequencies, namely $\pi$.
    \item $f$ is a reward per state: long-run average reward equals the $\pi$-weighted reward. With $f=\ind_{\{i\}}$ this is the first statement.
  \end{itemize}
\end{frame}

% ---------------------------------------------------------------------------
\begin{frame}{The Ergodic Theorem: Proof Idea and a Time Average}
  \label{fr:ergodicwhy}
  \small
  \begin{block}{Proof idea}
    The visits to $i$ split the path into independent, identically distributed excursions of mean length $\E_i[T_i]$ (Markov property at each return). By the law of large numbers (LLN: averages of i.i.d.\ variables with finite mean converge to that mean with probability one), $m$ excursions take about $m\,\E_i[T_i]$ steps, so in $n$ steps there are about $n/\E_i[T_i]$ visits: $V_i(n)/n\to\mu_i:=1/\E_i[T_i]$.
  \end{block}
  \begin{block}{Why the limit is $\pi$}
    Among the visits to $i$, a fraction $p(i,j)$ is followed by a step to $j$ (Markov property, LLN), so the fraction of times $k<n$ with $X_k=i$, $X_{k+1}=j$ tends to $\mu_i\,p(i,j)$. Summing over $i$ counts the visits to $j$ at times $1,\dots,n$, whose fraction tends to $\mu_j$: $\mu_j=\sum_i\mu_i\,p(i,j)$, i.e.\ $\mu\PP=\mu$, so $\mu=\pi$ by uniqueness (slide~\ref{fr:recclass}). \hfill$\square$
  \end{block}
  \words{The frequencies satisfy $\mu\PP=\mu$ because every arrival at $j$ is a departure from some $i$.}
  % python3: cube, f = number of 1s: sum_i pi_i f(i) = (0+3*1+3*2+3)/8 = 12/8 = 3/2; even-parity states: 4/8 = 1/2
  \begin{exampleblock}{The cube: a time average although the rows of $\PP^n$ never converge}
    $f(x)=$ number of $1$s in $x$. Along any path, with probability one, $\frac1n\sum_{k<n}f(X_k)\to\sum_i\pi_if(i)=\tfrac18(0+3\cdot1+3\cdot2+3)=\tfrac32$, and the fraction of time at even-parity corners tends to $\tfrac12$, although at each fixed time $n$ that probability is $1$ or $0$ (slide~\ref{fr:cubeparity}).
  \end{exampleblock}
\end{frame}

% ---------------------------------------------------------------------------
\begin{frame}{Return to Question 2, and the Two Theorems Side by Side}
  \label{fr:sidebyside}
  \small
  \begin{exampleblock}{Weather, question 2}
    In the long run a fraction $\pi_{\st S}=\tfrac6{13}\approx0.46$ of days are sunny, $\tfrac4{13}$ cloudy, $\tfrac3{13}$ rainy: with probability one, along the actual sequence of days. And the fraction of days that are rainy \emph{and} followed by a sunny day is $\pi_{\st R}\,p(\st R,\st S)=\tfrac3{13}\cdot0.2=\tfrac3{65}\approx0.046$ (apply the theorem to the pair chain $(X_n,X_{n+1})$).
    % python3: 3/13*0.2 = 3/65 = 0.04615
  \end{exampleblock}
  \begin{center}
    \renewcommand{\arraystretch}{1.4}
    \begin{tabular}{@{}lll@{}}
      \toprule
      & convergence theorem & ergodic theorem \\
      \midrule
      statement & $\Prob(X_n=i)\to\pi_i$ & $V_i(n)/n\to\pi_i$ with probability one \\
      reads the table & down a column (late time) & along a row (one path) \\
      needs & irreducible, aperiodic, finite & irreducible, finite \\
      the cube ($d=2$) & does not apply: rows of $\PP^n$ alternate & applies: $\tfrac18$ of the time at each corner \\
      what $\pi_i$ means & distribution of $X_n$ for large $n$ & long-run fraction of time at $i$ \\
      computes & $\PP^n$ for large $n$ & $\pi$ from $\pi\PP=\pi$, or $1/\E_i[T_i]$ \\
      \bottomrule
    \end{tabular}
  \end{center}
  \words{Both theorems describe the long-run behaviour of the chain in terms of $\pi$: the first as the distribution at a fixed late time, the second as the frequencies along a single path.}
\end{frame}


% ===========================================================================
\begin{frame}{Absorbing Chains and the Fundamental Matrix}
  \label{fr:absorbing}
  \small
  \begin{block}{Definition}
    A chain is \textbf{absorbing} if every state leads to some absorbing state. Order the states transient first, absorbing last: with $S_T$ the transient and $S_A$ the absorbing states,
    \[
      \PP=\begin{pmatrix}\mathbf Q&\mathbf R\\ \mathbf 0&\mathbf I\end{pmatrix},\qquad
      \mathbf Q\colon S_T\times S_T\to[0,1],\quad \mathbf R\colon S_T\times S_A\to[0,1].
    \]
    The \textbf{fundamental matrix} is $\mathbf N=(\mathbf I-\mathbf Q)^{-1}\colon S_T\times S_T\to[0,\infty)$.
  \end{block}
  \begin{block}{Theorem}
    $\mathbf I-\mathbf Q$ is invertible and $\mathbf N=\sum_{n\ge0}\mathbf Q^n$. For transient $i,j$ and absorbing $a$:
    \[
      \mathbf N_{ij}=\E_i\big[\#\{n\ge0 : X_n=j\}\big],\qquad
      (\mathbf N\mathbf 1)_i=\E_i[T_{S_A}],\qquad
      (\mathbf N\mathbf R)_{ia}=\Prob_i(\text{absorbed at } a).
    \]
  \end{block}
  \words{$\mathbf Q^n$ contains the $n$-step transition probabilities among the transient states, so its sum over $n$ gives expected numbers of visits; the row sums of $\mathbf N$ are expected times to absorption, and $\mathbf N\mathbf R$ gives the absorption probabilities.}
\end{frame}

% ---------------------------------------------------------------------------
\begin{frame}{The Fundamental Matrix: Proof Sketch}
  \label{fr:fundproof}
  \small
  \begin{block}{Proof sketch}
    $(\mathbf Q^n)_{ij}=\Prob_i(X_n=j)$ for transient $j$ (paths through transient states only), so $\sum_n(\mathbf Q^n)_{ij}=\E_i[\text{visits to } j]$ by linearity, column by column. Absorption is certain and from a finite set of transient states happens geometrically fast, so $\mathbf Q^n\to0$ and the series converges, to $(\mathbf I-\mathbf Q)^{-1}$ since $(\mathbf I-\mathbf Q)\sum_{n<m}\mathbf Q^n=\mathbf I-\mathbf Q^m$. Time to absorption is the total number of transient visits: $\sum_j\mathbf N_{ij}$. Absorption at $a$ occurs by a visit to some transient $j$ followed by a step to $a$, giving $\sum_j\mathbf N_{ij}\mathbf R_{ja}$. \hfill$\square$
  \end{block}
  \words{The expected numbers of visits, the expected time to absorption and the absorption probabilities are all obtained from the single matrix $(\mathbf I-\mathbf Q)^{-1}$.}
  \begin{itemize}
    \item The hitting systems of slides~\ref{fr:hitting} and~\ref{fr:khit} are $(\mathbf I-\mathbf Q)\,k=\mathbf 1$ and $(\mathbf I-\mathbf Q)\,h=\mathbf R_{\cdot,a}$; the fundamental matrix is their solution operator.
    \item For a hitting question, any chain can be made absorbing by making the target states absorbing (replacing their rows by point masses), which does not change the path up to the hitting time.
  \end{itemize}
\end{frame}

% ---------------------------------------------------------------------------
\begin{frame}{The Fundamental Matrix of the Ruin Chain, $N=4$}
  \label{fr:fundruin}
  \small
  % python3 sympy: p=1/2: Q=[[0,1/2,0],[1/2,0,1/2],[0,1/2,0]], R=[[1/2,0],[0,0],[0,1/2]];
  %   N=(I-Q)^-1=[[3/2,1,1/2],[1,2,1],[1/2,1,3/2]]; N1=(3,4,3); NR=[[3/4,1/4],[1/2,1/2],[1/4,3/4]]
  %   p=2/5: N=[[19/13,10/13,4/13],[15/13,25/13,10/13],[9/13,15/13,19/13]]; N1=(33/13,50/13,43/13);
  %   NR=[[57/65,8/65],[9/13,4/13],[27/65,38/65]]; matches win(i,4,2/5)=(8/65,4/13,38/65) and dur=(33/13,50/13,43/13).
  Transient states $1,2,3$; absorbing $0,4$. Fair case:
  \[
    \mathbf Q=\begin{pmatrix}0&\tfrac12&0\\ \tfrac12&0&\tfrac12\\ 0&\tfrac12&0\end{pmatrix},\quad
    \mathbf R=\begin{pmatrix}\tfrac12&0\\0&0\\0&\tfrac12\end{pmatrix},\quad
    \mathbf N=(\mathbf I-\mathbf Q)^{-1}=\begin{pmatrix}\tfrac32&1&\tfrac12\\1&2&1\\ \tfrac12&1&\tfrac32\end{pmatrix},
  \]
  \[
    \mathbf N\mathbf 1=\begin{pmatrix}3\\4\\3\end{pmatrix},\qquad
    \mathbf N\mathbf R=\begin{pmatrix}\tfrac34&\tfrac14\\ \tfrac12&\tfrac12\\ \tfrac14&\tfrac34\end{pmatrix}.
  \]
  \begin{exampleblock}{Reading it}
    $\mathbf N\mathbf 1=(3,4,3)=i(4-i)$: the durations of slide~\ref{fr:duration}. Column $4$ of $\mathbf N\mathbf R$ is $(\tfrac14,\tfrac12,\tfrac34)=i/4$: slide~\ref{fr:ruinfair}. The matrix also gives new information: $\mathbf N_{2,2}=2$ is the expected number of visits to fortune $2$ (counting the start) before the game ends when starting there, and $\mathbf N_{1,3}=\tfrac12$ the expected number of visits to $3$ when starting from $1$.
  \end{exampleblock}
  \begin{itemize}
    \item $p=\tfrac25$: $\mathbf N\mathbf 1=(\tfrac{33}{13},\tfrac{50}{13},\tfrac{43}{13})$ and the target column of $\mathbf N\mathbf R$ is $(\tfrac8{65},\tfrac4{13},\tfrac{38}{65})$, agreeing with the closed forms of slides~\ref{fr:ruinunfair}--\ref{fr:duration}. The single matrix inverse gives both quantities at once.
  \end{itemize}
\end{frame}


% ===========================================================================
\begin{frame}{Summary}
  \label{fr:summary}
  \footnotesize
  \renewcommand{\arraystretch}{1.15}
  \begin{tabular}{@{}>{\raggedright\arraybackslash}p{3.2cm}>{\raggedright\arraybackslash}p{3.7cm}>{\raggedright\arraybackslash}p{6.1cm}@{}}
    \toprule
    object & type & meaning in the table of paths \\
    \midrule
    process $(X_t)_{t\in T}$ & $X_t\colon\Omega\to\R$ for each $t$ & columns are random variables, rows are paths \\
    random walk $S_n$ & $\Omega\to\Z$ & sum of $\pm1$ steps; $\Var(S_n)=4npq$ ($=n$ if fair); restarts after each step \\
    Brownian motion $B_t$ & $[0,\infty)\times\Omega\to\R$ & the walk at $\sqrt n$ scale; independent $N(0,t-s)$ increments \\
    ruin probability $1-h_i$ & $\{0,\dots,N\}\to[0,1]$ & $1-i/N$, or $\frac{r^N-r^i}{r^N-1}$; from $h_i=ph_{i+1}+qh_{i-1}$ \\
    transition matrix $\PP$ & $S\times S\to[0,1]$, rows sum to $1$ & row $i$ is the conditional pmf of $X_{n+1}$ given $X_n=i$ \\
    $n$-step $\PP^n$, law $\lambda\PP^n$ & matrix, row vector & Chapman--Kolmogorov: condition on the middle state \\
    stationary $\pi$ & $S\to[0,1]$, $\pi\PP=\pi$ & inflow equals outflow at every state; detailed balance if reversible \\
    classes, period, recurrence & structure of $\PP$'s graph & which states reach which, at what step lengths, and whether return is certain \\
    hitting $h^A$, $k^A$ & $S\to[0,1]$, $S\to[0,\infty]$ & minimal solutions of first-step systems; $\mathbf N\mathbf R$, $\mathbf N\mathbf 1$ \\
    convergence, ergodic & $\PP^n\to\mathbf 1\pi$; $V_i(n)/n\to\pi_i$ & the column limit and the row limit are both $\pi$ \\
    \bottomrule
  \end{tabular}
  \medskip
  \words{Most results follow from conditioning on the first step, turning a question about unbounded paths into a finite linear system.}
\end{frame}
